
\subsection{假先知/假使徒/尼哥拉一黨人/假弟兄/猶太人/撒旦一會}
\subsubsection{假先知}
 耶利米書 23:9-40\\
論到那些先知，我心在我裡面憂傷，我骨頭都發顫。因耶和華和他的聖言，我像醉酒的人，像被酒所勝的人。 10 地滿了行淫的人。因妄自賭咒，地就悲哀，曠野的草場都枯乾了。他們所行的道乃是惡的，他們的勇力使得不正。 11 「連先知帶祭司都是褻瀆的，就是在我殿中我也看見他們的惡。」這是耶和華說的。 12 「因此，他們的道路必像黑暗中的滑地，他們必被追趕在這路中仆倒。因為當追討之年，我必使災禍臨到他們。」這是耶和華說的。13 「我在撒馬利亞的先知中曾見愚妄，他們藉巴力說預言，使我的百姓以色列走錯了路。 14 我在耶路撒冷的先知中曾見可憎惡的事，他們行姦淫，做事虛妄，又堅固惡人的手，甚至無人回頭離開他的惡。他們在我面前都像所多瑪，耶路撒冷的居民都像蛾摩拉。」 15 所以萬軍之耶和華論到先知如此說：「我必將茵陳給他們吃，又將苦膽水給他們喝，因為褻瀆的事出於耶路撒冷的先知，流行遍地。」16 萬軍之耶和華如此說：「這些先知向你們說預言，你們不要聽他們的話。他們以虛空教訓你們，所說的異象是出於自己的心，不是出於耶和華的口。 17 他們常對藐視我的人說：『耶和華說，你們必享平安。』又對一切按自己頑梗之心而行的人說：『必沒有災禍臨到你們。』」 18 有誰站在耶和華的會中得以聽見並會悟他的話呢？有誰留心聽他的話呢？ 19 看哪，耶和華的憤怒好像暴風，已經發出，是暴烈的旋風，必轉到惡人的頭上。 20 耶和華的怒氣必不轉消，直到他心中所擬定的成就了，末後的日子你們要全然明白。 21 「我沒有打發那些先知，他們竟自奔跑；我沒有對他們說話，他們竟自預言。 22 他們若是站在我的會中，就必使我的百姓聽我的話，又使他們回頭離開惡道和他們所行的惡。」23 耶和華說：「我豈為近處的神呢？不也為遠處的神嗎？」 24 耶和華說：「人豈能在隱密處藏身，使我看不見他呢？」耶和華說：「我豈不充滿天地嗎？ 25 我已聽見那些先知所說的，就是託我名說的假預言，他們說：『我做了夢！我做了夢！』 26 說假預言的先知，就是預言本心詭詐的先知，他們這樣存心要到幾時呢？ 27 他們各人將所做的夢對鄰舍述說，想要使我的百姓忘記我的名，正如他們列祖因巴力忘記我的名一樣。 28 得夢的先知可以述說那夢，得我話的人可以誠實講說我的話！糠秕怎能與麥子比較呢？」這是耶和華說的。 29 耶和華說：「我的話豈不像火，又像能打碎磐石的大錘嗎？」 30 耶和華說：「那些先知各從鄰舍偷竊我的言語，因此我必與他們反對。」 31 耶和華說：「那些先知用舌頭說『是耶和華說的』，我必與他們反對。」 32 耶和華說：「那些以幻夢為預言，又述說這夢，以謊言和矜誇使我百姓走錯了路的，我必與他們反對。我沒有打發他們，也沒有吩咐他們，他們與這百姓毫無益處。」這是耶和華說的。
33 「無論是百姓，是先知，是祭司，問你說：『耶和華有什麼默示呢？』你就對他們說：『什麼默示啊？耶和華說：我要撇棄你們！』 34 無論是先知，是祭司，是百姓，說『耶和華的默示』，我必刑罰那人和他的家。 35 你們各人要對鄰舍，各人要對弟兄如此說：『耶和華回答什麼？耶和華說了什麼呢？』 36 『耶和華的默示』你們不可再提，各人所說的話必做自己的重擔 a，因為你們謬用永生神萬軍之耶和華我們神的言語。 37 你們要對先知如此說：『耶和華回答你什麼？耶和華說了什麼呢？』 38 你們若說『耶和華的默示』，耶和華就如此說：因你們說『耶和華的默示』這句話，我也打發人到你們那裡去，告訴你們不可說『耶和華的默示』， 39 所以我必全然忘記你們，將你們和我所賜給你們並你們列祖的城撇棄了。 40 又必使永遠的凌辱和長久的羞恥臨到你們，是不能忘記的。」

耶利米書 27:16-18\\
我又對祭司和這眾民說：「耶和華如此說：你們不可聽那先知對你們所說的預言，他們說『耶和華殿中的器皿快要從巴比倫帶回來』，其實他們向你們說假預言。 17 不可聽從他們，只管服侍巴比倫王便得存活。這城何至變為荒場呢？ 18 他們若果是先知，有耶和華的話臨到他們，讓他們祈求萬軍之耶和華，使那在耶和華殿中和猶大王宮內，並耶路撒冷剩下的器皿不被帶到巴比倫去。

使徒行傳 20：25-35 \\
 「我素常在你們中間來往，傳講神國的道，如今我曉得，你們以後都不得再見我的面了。 26 所以我今日向你們證明，你們中間無論何人死亡，罪不在我身上 c， 27 因為神的旨意我並沒有一樣避諱不傳給你們的。 28 聖靈立你們做全群的監督，你們就當為自己謹慎，也為全群謹慎，牧養神的教會，就是他用自己血所買來的 d。 29 我知道我去之後，必有凶暴的豺狼進入你們中間，不愛惜羊群。 30 就是你們中間也必有人起來，說悖謬的話，要引誘門徒跟從他們。 31 所以你們應當警醒，記念我三年之久晝夜不住地流淚勸誡你們各人。 32 如今我把你們交託神和他恩惠的道，這道能建立你們，叫你們和一切成聖的人同得基業。 33 我未曾貪圖一個人的金、銀、衣服。 34 我這兩隻手常供給我和同人的需用，這是你們自己知道的。 35 我凡事給你們做榜樣，叫你們知道應當這樣勞苦扶助軟弱的人，又當記念主耶穌的話說：『施比受更為有福。』」
 
 加拉太書 1
 我稀奇你們這麼快離開那藉著基督之恩召你們的，去從別的福音。 7 那並不是福音，不過有些人攪擾你們，要把基督的福音更改了。 8 但無論是我們，是天上來的使者，若傳福音給你們，與我們所傳給你們的不同，他就應當被咒詛。 9 我們已經說了，現在又說：若有人傳福音給你們，與你們所領受的不同，他就應當被咒詛！ 10 我現在是要得人的心呢，還是要得神的心呢？我豈是討人的喜歡嗎？若仍舊討人的喜歡，我就不是基督的僕人了。
 
 \subsubsection{假使徒}
 \subsubsubsection{以弗所 }
 我知道你的行為、勞碌、忍耐，也知道你不能容忍惡人。你也曾試驗那自稱為使徒卻不是使徒的，看出他們是假的來。
 
 \subsubsubsection{哥林多後書 11：13-23}
 13 那等人是假使徒，行事詭詐，裝作基督使徒的模樣。 14 這也不足為怪，因為連撒旦也裝作光明的天使。 15 所以，他的差役若裝作仁義的差役，也不算稀奇。他們的結局必然照著他們的行為。16 我再說，人不可把我看做愚妄的；縱然如此，也要把我當做愚妄人接納，叫我可以略略自誇。 17 我說的話不是奉主命說的，乃是像愚妄人放膽自誇； 18 既有好些人憑著血氣自誇，我也要自誇了。 19 你們既是精明人，就能甘心忍耐愚妄人！ 20 假若有人強你們做奴僕，或侵吞你們，或擄掠你們，或侮慢你們，或打你們的臉，你們都能忍耐他！ 21 我說這話是羞辱自己，好像我們從前是軟弱的！然而，人在何事上勇敢，我說句愚妄話，我也勇敢！ 22 他們是希伯來人嗎？我也是。他們是以色列人嗎？我也是。他們是亞伯拉罕的後裔嗎？我也是。 23 他們是基督的僕人嗎？我說句狂話，我更是！
 
 \subsubsection{尼哥拉一黨人——以弗所，別迦摩}
 \begin{itemize}
 \itemsep-.25em
\item 然而，你還有一件可取的事，就是你恨惡尼哥拉一黨人的行為，這也是我所恨惡的。
\item 你那裡(別迦摩)也有人照樣服從了尼哥拉一黨人的教訓。
 \end{itemize} 
 
 \subsubsection{撒旦一會——士每拿,非拉鐵非}
 \begin{itemize}
 \itemsep-.25em
\item  「你要寫信給士每拿教會的使者說：『那首先的、末後的、死過又活的說： 9 我知道你的患難、你的貧窮，你卻是富足的；也知道那自稱是猶太人所說的毀謗話，其實他們不是猶太人，乃是撒旦一會的人。
\item  那撒旦一會的，自稱是猶太人，其實不是猶太人，乃是說謊話的，我要使他們來在你腳前下拜，也使他們知道我是已經愛你了。
 \end{itemize}
 
 
 \subsubsection{假弟兄——加拉太書（2：4-8），律法（3：1-4），割礼（5：5-12，6：12-12 ） } 
 \begin{itemize}
 \itemsep-.25em
\item 加拉太書 2：4-8
  4 因為有偷著引進來的假弟兄，私下窺探我們在基督耶穌裡的自由，要叫我們做奴僕。 5 我們就是一刻的工夫也沒有容讓順服他們，為要叫福音的真理仍存在你們中間。 6 至於那些有名望的，不論他是何等人，都與我無干，神不以外貌取人。那些有名望的並沒有加增我什麼， 7 反倒看見了主託我傳福音給那未受割禮的人，正如託彼得傳福音給那受割禮的人。 8 那感動彼得叫他為受割禮之人做使徒的，也感動我，叫我為外邦人做使徒。
\item  加拉太書 3：1-4
 無知的加拉太人哪！耶穌基督釘十字架，已經活畫在你們眼前，誰又迷惑了你們呢？ 2 我只要問你們這一件：你們受了聖靈，是因行律法呢，是因聽信福音呢？ 3 你們既靠聖靈入門，如今還靠肉身成全嗎？你們是這樣的無知嗎？ 4 你們受苦如此之多，都是徒然的嗎？難道果真是徒然的嗎？
 \item 加拉太書 5：5-12
  5 我們靠著聖靈，憑著信心，等候所盼望的義。 6 原來在基督耶穌裡，受割禮不受割禮全無功效，唯獨使人生發仁愛的信心才有功效。 7 你們向來跑得好，有誰攔阻你們，叫你們不順從真理呢？ 8 這樣的勸導不是出於那召你們的。 9 一點麵酵能使全團都發起來。 10 我在主裡很信你們必不懷別樣的心，但攪擾你們的，無論是誰，必擔當他的罪名。 11 弟兄們，我若仍舊傳割禮，為什麼還受逼迫呢？若是這樣，那十字架討厭的地方就沒有了。 12 恨不得那攪亂你們的人把自己割絕了！
\item 加拉太書 6：12-13 
   12 凡希圖外貌體面的人，都勉強你們受割禮，無非是怕自己為基督的十字架受逼迫。 13 他們那些受割禮的，連自己也不守律法，他們願意你們受割禮，不過要藉著你們的肉體誇口。
  \end{itemize}
 
  \subsubsubsection{腓立比書  割礼（3:2-3,17-19） } 
  \begin{itemize}
 \itemsep-.25em
\item  2 應當防備犬類，防備作惡的，防備妄自行割的。 3 因為真受割禮的，乃是我們這以神的靈敬拜，在基督耶穌裡誇口，不靠著肉體的。
 \item 17 弟兄們，你們要一同效法我，也當留意看那些照我們榜樣行的人。 18 因為有許多人行事是基督十字架的仇敵，我屢次告訴你們，現在又流淚地告訴你們。 19 他們的結局就是沉淪，他們的神就是自己的肚腹！他們以自己的羞辱為榮耀，專以地上的事為念。
 \end{itemize}
 
 \subsubsubsection{提前 （1:3-11,18-20,3:9-11,4:1-5,6:3-10） } 
  \begin{itemize}
 \itemsep-.25em
\item 3 我往馬其頓去的時候，曾勸你仍住在以弗所，好囑咐那幾個人不可傳異教， 4 也不可聽從荒渺無憑的話語和無窮的家譜；這等事只生辯論，並不發明神在信上所立的章程。 5 但命令的總歸就是愛，這愛是從清潔的心和無虧的良心、無偽的信心生出來的。 6 有人偏離這些，反去講虛浮的話， 7 想要做教法師，卻不明白自己所講說的、所論定的。 8 我們知道律法原是好的，只要人用得合宜。 9 因為律法不是為義人設立的，乃是為不法和不服的，不虔誠和犯罪的，不聖潔和戀世俗的，弑父母和殺人的， 10 行淫和親男色的，搶人口和說謊話的，並起假誓的，或是為別樣敵正道的事設立的。 11 這是照著可稱頌之神交託我榮耀福音說的。
 \item  我兒提摩太啊，我照從前指著你的預言，將這命令交託你，叫你因此可以打那美好的仗， 19 常存信心和無虧的良心。有人丟棄良心，就在真道上如同船破壞了一般。 20 其中有許米乃和亞歷山大，我已經把他們交給撒旦，使他們受責罰就不再謗讟了。
  \item 聖靈明說，在後來的時候，必有人離棄真道，聽從那引誘人的邪靈和鬼魔的道理。 2 這是因為說謊之人的假冒，這等人的良心如同被熱鐵烙慣了一般。 3 他們禁止嫁娶，又禁戒食物 a，就是神所造，叫那信而明白真道的人感謝著領受的。 4 凡神所造的物都是好的，若感謝著領受，就沒有一樣可棄的， 5 都因神的道和人的祈求成為聖潔了。
 \item  3 若有人傳異教，不服從我們主耶穌基督純正的話與那合乎敬虔的道理， 4 他是自高自大，一無所知，專好問難，爭辯言辭，從此就生出嫉妒、紛爭、毀謗、妄疑， 5 並那壞了心術、失喪真理之人的爭競。他們以敬虔為得利的門路。 6 然而，敬虔加上知足的心便是大利了。 7 因為我們沒有帶什麼到世上來，也不能帶什麼去。 8 只要有衣有食，就當知足。9 但那些想要發財的人，就陷在迷惑，落在網羅和許多無知有害的私慾裡，叫人沉在敗壞和滅亡中。 10 貪財是萬惡之根。有人貪戀錢財，就被引誘離了真道，用許多愁苦把自己刺透了。
 \end{itemize}
    
\subsubsection{提多書  猶太人（1:10-19,3:9-11） } 
  \begin{itemize}
 \itemsep-.25em
\item  10 因為有許多人不服約束，說虛空話欺哄人，那奉割禮的更是這樣。 11 這些人的口總要堵住。他們因貪不義之財，將不該教導的教導人，敗壞人的全家。 12 有克里特人中的一個本地先知說：「克里特人常說謊話，乃是惡獸，又饞又懶。」 13 這個見證是真的。所以，你要嚴嚴地責備他們，使他們在真道上純全無疵， 14 不聽猶太人荒渺的言語和離棄真道之人的誡命。15 在潔淨的人，凡物都潔淨；在汙穢不信的人，什麼都不潔淨，連心地和天良也都汙穢了。 16 他們說是認識神，行事卻和他相背；本是可憎惡的，是悖逆的，在各樣善事上是可廢棄的。
 \item  9 要遠避無知的辯論和家譜的空談，以及紛爭並因律法而起的爭競，因為這都是虛妄無益的。 10 分門結黨的人，警戒過一兩次，就要棄絕他。 11 因為知道這等人已經背道，犯了罪，自己明知不是，還是去做。
 \end{itemize}  
    
 \subsubsection{巴蘭的教訓——別迦摩}
 然而，有幾件事我要責備你，因為在你那裡，有人服從了巴蘭的教訓。這巴蘭曾教導巴勒將絆腳石放在以色列人面前，叫他們吃祭偶像之物，行姦淫的事。
 
 \subsubsection{耶洗別——推雅推喇}
 然而，有一件事我要責備你，就是你容讓那自稱是先知的婦人耶洗別教導我的僕人，引誘他們行姦淫，吃祭偶像之物。 21 我曾給她悔改的機會，她卻不肯悔改她的淫行。 22 看哪，我要叫她病臥在床；那些與她行淫的人，若不悔改所行的，我也要叫他們同受大患難。 23 我又要殺死她的黨類 a，叫眾教會知道我是那察看人肺腑心腸的，並要照你們的行為報應你們各人。
 
 \subsubsection{撒旦深奧之理——推雅推喇} 
 『至於你們推雅推喇其餘的人，就是一切不從那教訓、不曉得他們素常所說撒旦深奧之理的人，我告訴你們：我不將別的擔子放在你們身上； 25 但你們已經有的，總要持守，直等到我來。
 
 \subsubsection{敵基督-約翰一書 2:18-23????????????} 
18 小子們哪，如今是末時了。你們曾聽見說那敵基督的要來，現在已經有好些敵基督的出來了；從此，我們就知道如今是末時了。 19 他們從我們中間出去，卻不是屬我們的，若是屬我們的，就必仍舊與我們同在；他們出去，顯明都不是屬我們的。 20 你們從那聖者受了恩膏，並且知道這一切的事 a。 21 我寫信給你們，不是因你們不知道真理，正是因你們知道，並且知道沒有虛謊是從真理出來的。 22 誰是說謊話的呢？不是那不認耶穌為基督的嗎？不認父與子的，這就是敵基督的。 23 凡不認子的，就沒有父；認子的，連父也有了。







\subsection{金燈臺和二橄欖樹指二受膏者}

\subsubsection{啟示錄11：3-13}
 3 我要使我那兩個見證人穿著毛衣，傳道一千二百六十天。」 4 他們就是那兩棵橄欖樹，兩個燈臺，立在世界之主面前的。 5 若有人想要害他們，就有火從他們口中出來，燒滅仇敵。凡想要害他們的都必這樣被殺。 6 這二人有權柄在他們傳道的日子叫天閉塞不下雨，又有權柄叫水變為血，並且能隨時隨意用各樣的災殃攻擊世界。 7 他們作完見證的時候，那從無底坑裡上來的獸必與他們交戰，並且得勝，把他們殺了。8 他們的屍首就倒在大城裡的街上，這城按著靈意叫所多瑪，又叫埃及，就是他們的主釘十字架之處。 9 從各民、各族、各方、各國中，有人觀看他們的屍首三天半，又不許把屍首放在墳墓裡。 10 住在地上的人就為他們歡喜快樂，互相饋送禮物，因這兩位先知曾叫住在地上的人受痛苦。過了這三天半，有生氣從神那裡進入他們裡面，他們就站起來，看見他們的人甚是害怕。 12 兩位先知聽見有大聲音從天上來，對他們說：「上到這裡來！」他們就駕著雲上了天，他們的仇敵也看見了。 13 正在那時候，地大震動，城就倒塌了十分之一，因地震而死的有七千人，其餘的都恐懼，歸榮耀給天上的神。

\subsubsection{撒迦利亞書 4：1-14}
那與我說話的天使又來叫醒我，好像人睡覺被喚醒一樣。 2 他問我說：「你看見了什麼？」我說：「我看見了一個純金的燈臺，頂上有燈盞，燈臺上有七盞燈，每盞有七個管子。 3 旁邊有兩棵橄欖樹，一棵在燈盞的右邊，一棵在燈盞的左邊。」 4 我問與我說話的天使說：「主啊，這是什麼意思？」 5 與我說話的天使回答我說：「你不知道這是什麼意思嗎？」我說：「主啊，我不知道。」 6 他對我說：「這是耶和華指示所羅巴伯的。萬軍之耶和華說：不是倚靠勢力，不是倚靠才能，乃是倚靠我的靈方能成事。 7 大山哪，你算什麼呢？在所羅巴伯面前，你必成為平地。他必搬出一塊石頭，安在殿頂上，人且大聲歡呼說：『願恩惠，恩惠歸於這殿 a！』」 8 耶和華的話又臨到我說： 9 「所羅巴伯的手立了這殿的根基，他的手也必完成這工，你就知道萬軍之耶和華差遣我到你們這裡來了。 10 誰藐視這日的事為小呢？這七眼乃是耶和華的眼睛，遍察全地，見所羅巴伯手拿線砣就歡喜。」
我又問天使說：「這燈臺左右的兩棵橄欖樹是什麼意思？」 12 我二次問他說：「這兩根橄欖枝在兩個流出金色油的金嘴旁邊是什麼意思？」 13 他對我說：「你不知道這是什麼意思嗎？」我說：「主啊，我不知道。」 14 他說：「這是兩個受膏者，站在普天下主的旁邊。」






\subsection{谨慎}
\subsubsection{谨慎}
%馬太福音 18:10\\
%這世界有禍了，因為將人絆倒；絆倒人的事是免不了的，但那絆倒人的有禍了！ 8 倘若你一隻手或是一隻腳叫你跌倒，就砍下來丟掉！你缺一隻手或是一隻腳進入永生，強如有兩手兩腳被丟在永火裡。 9 倘若你一隻眼叫你跌倒，就把它剜出來丟掉！你只有一隻眼進入永生，強如有兩隻眼被丟在地獄的火裡。 10 你們要\textcolor{red}{小心}，不可輕看這小子裡的一個。我告訴你們，他們的使者在天上，常見我天父的面。


馬太福音 24:7-14\\
耶穌回答說：「你們要\textcolor{red}{谨慎}，免得有人迷惑你們。 5 因為將來有好些人冒我的名來，說：『我是基督』，並且要迷惑許多人。 6 你們也要聽見打仗和打仗的風聲，總不要驚慌，因為這些事是必須有的，\textcolor{blue}{只是末期還沒有到}。
民要攻打民，國要攻打國，多處必有饑荒、地震， 8 這都是災難 a的起頭。 9 那時，人要把你們陷在患難裡，也要殺害你們；你們又要為我的名被萬民恨惡。10 「那時，必有許多人跌倒，也要彼此陷害、彼此恨惡， 11 且有好些假先知起來，迷惑多人。 12 只因不法的事增多，許多人的愛心才漸漸冷淡了。 13 唯有忍耐到底的，必然得救。 14 這天國的福音要傳遍天下，對萬民作見證，\textcolor{blue}{然後末期才來到}。

马可福音 4:24\\
你們所聽的要\textcolor{red}{留心}。你們用什麼量器量給人，也必用什麼量器量給你們，並且要多給你們。 25 因為有的，還要給他；沒有的，連他所有的也要奪去。


马可福音 8:14-15\\
門徒忘了帶餅，在船上除了一個餅，沒有別的食物。 15 耶穌囑咐他們說：「你們要\textcolor{red}{谨慎}，防備法利賽人的酵和希律的酵。」

马可福音 12:38-40\\
耶穌在教訓之間說：「你們要\textcolor{red}{防備}文士。他們好穿長衣遊行，喜愛人在街市上問他們的安， 39 又喜愛會堂裡的高位、筵席上的首座。 40 他們侵吞寡婦的家產，假意作很長的禱告。這些人要受更重的刑罰！」

马可福音 13:5-13\\
耶稣说：“你们要\textcolor{red}{谨慎}，免得有人迷惑你们。 6 将来有好些人冒我的名来，说‘我是基督’，并且要迷惑许多人。 7 你们听见打仗和打仗的风声，不要惊慌，这些事是必须有的，只是末期还没有到。 8 民要攻打民，国要攻打国，多处必有地震、饥荒，这都是灾难的起头。但你们要\textcolor{red}{谨慎}，因为人要把你们交给公会，并且你们在会堂里要受鞭打，又为我的缘故站在诸侯与君王面前，对他们作见证。 10 然而，福音必须先传给万民。 11 人把你们拉去交官的时候，不要预先思虑说什么，到那时候，赐给你们什么话，你们就说什么；因为说话的不是你们，乃是圣灵。 12 弟兄要把弟兄，父亲要把儿子，送到死地；儿女要起来与父母为敌，害死他们。 13 并且你们要为我的名被众人恨恶，唯有忍耐到底的，必然得救。

马可福音 13:21-23\\
那時，若有人對你們說『看哪，基督在這裡』，或說『基督在那裡』，你們不要信。 22 因為假基督、假先知將要起來，顯神蹟奇事，倘若能行，就把選民迷惑了。 23 你們要\textcolor{red}{謹慎}！看哪，凡事我都預先告訴你們了！

马可福音 13:33-36\\
你们要\textcolor{red}{谨慎}，警醒祈祷，因为你们不晓得那日期几时来到。 34 这事正如一个人离开本家，寄居外邦，把权柄交给仆人，分派各人当做的工，又吩咐看门的警醒。 35 所以，你们要警醒，因为你们不知道家主什么时候来，或晚上，或半夜，或鸡叫，或早晨。 36 恐怕他忽然来到，看见你们睡着了。 37 我对你们所说的话，也是对众人说：要警醒！

路加福音 8:16-18\\
沒有人點燈用器皿蓋上或放在床底下，乃是放在燈臺上，叫進來的人看見亮光。 17 因為掩藏的事沒有不顯出來的，隱瞞的事沒有不露出來被人知道的。 18 所以，你們應當\textcolor{red}{小心}怎樣聽。因為凡有的，還要加給他；凡沒有的，連他自以為有的也要奪去。

路加福音 21:8-9\\
耶穌說：「你們要\textcolor{red}{謹慎}，不要受迷惑。因為將來有好些人冒我的名來，說『我是基督』，又說『時候近了』。你們不要跟從他們！ 9 你們聽見打仗和擾亂的事，不要驚惶，因為這些事必須先有，只是末期不能立時就到。」

路加福音 21:34-36\\
你们要\textcolor{red}{謹慎}，恐怕因贪食、醉酒并今生的思虑累住你们的心，那日子就如同网罗忽然临到你们， 35 因为那日子要这样临到全地上一切居住的人。 36 你们要时时警醒，常常祈求，使你们能逃避这一切要来的事，得以站立在人子面前。

使徒行傳 13:38-41\\
所以弟兄們，你們當曉得：赦罪的道是由這人傳給你們的！ 39 你們靠摩西的律法，在一切不得稱義的事上信靠這人，就都得稱義了。 40 所以你們務要\textcolor{red}{小心}，免得先知書上所說的臨到你們。 41 主說：『你們這輕慢的人，要觀看，要驚奇，要滅亡！因為在你們的時候，我行一件事，雖有人告訴你們，你們總是不信。』

使徒行傳 20:25-31\\
我素常在你們中間來往，傳講神國的道，如今我曉得，你們以後都不得再見我的面了。 26 所以我今日向你們證明，你們中間無論何人死亡，罪不在我身上， 27 因為神的旨意我並沒有一樣避諱不傳給你們的。 28 聖靈立你們做全群的監督，你們就當為自己\textcolor{red}{謹慎}，也為全群\textcolor{red}{謹慎}，牧養神的教會，就是他用自己血所買來的 d。 29 我知道我去之後，必有凶暴的豺狼進入你們中間，不愛惜羊群。 30 就是你們中間也必有人起來，說悖謬的話，要引誘門徒跟從他們。 31 所以你們應當警醒，記念我三年之久晝夜不住地流淚勸誡你們各人。

林前 3:10-15\\
我照神所給我的恩，好像一個聰明的工頭立好了根基，有別人在上面建造，只是各人要\textcolor{red}{謹慎}怎樣在上面建造。 11 因為那已經立好的根基就是耶穌基督，此外沒有人能立別的根基。 12 若有人用金、銀、寶石、草、木、禾秸在這根基上建造， 13 各人的工程必然顯露；因為那日子要將它表明出來，有火發現，這火要試驗各人的工程怎樣。 14 人在那根基上所建造的工程若存得住，他就要得賞賜； 15 人的工程若被燒了，他就要受虧損，自己卻要得救，雖然得救，乃像從火裡經過的一樣。

林前 10:11-12\\
他們遭遇這些事都要作為鑒戒，並且寫在經上正是警戒我們這末世的人。 12 所以，自己以為站得穩的須要\textcolor{red}{謹慎}，免得跌倒。

林前 16:10-11\\
若是提摩太來到，你們要\textcolor{red}{留心}，叫他在你們那裡無所懼怕，因為他勞力做主的工，像我一樣。 11 所以，無論誰都不可藐視他，只要送他平安前行，叫他到我這裡來，因我指望他和弟兄們同來。

加拉太書 5:13-15
弟兄們，你們蒙召是要得自由，只是不可將你們的自由當做放縱情慾的機會，總要用愛心互相服侍。 14 因為全律法都包在「愛人如己」這一句話之內了。 15 你們要\textcolor{red}{謹慎}，若相咬相吞，只怕要彼此消滅了。

腓立比書 3:2-3
應當\textcolor{red}{防備}犬類，\textcolor{red}{防備}作惡的，\textcolor{red}{防備}妄自行割的。 3 因為真受割禮的，乃是我們這以神的靈敬拜，在基督耶穌裡誇口，不靠著肉體的。

歌羅西書 2:8
你們要\textcolor{red}{謹慎}，恐怕有人用他的理學和虛空的妄言，不照著基督，乃照人間的遺傳和世上的小學，就把你們擄去。

希伯來書 3:12-13
弟兄們，你們要\textcolor{red}{謹慎}，免得你們中間或有人存著不信的惡心，把永生神離棄了。 13 總要趁著還有今日，天天彼此相勸，免得你們中間有人被罪迷惑，心裡就剛硬了。

希伯來書 12:22-25
你們乃是來到錫安山，永生神的城邑，就是天上的耶路撒冷；那裡有千萬的天使， 23 有名錄在天上諸長子之會所共聚的總會，有審判眾人的神和被成全之義人的靈魂， 24 並新約的中保耶穌以及所灑的血。這血所說的比亞伯的血所說的更美。 25 你們總要\textcolor{red}{謹慎}，不可棄絕那向你們說話的。因為那些棄絕在地上警戒他們的，尚且不能逃罪，何況我們違背那從天上警戒我們的呢？

彼得前書 4:7\\
萬物的結局近了，所以你們要你们要\textcolor{red}{謹慎}自守，警醒禱告。

彼得前書 5:8
務要你们要\textcolor{red}{謹守}、警醒，因為你們的仇敵魔鬼如同吼叫的獅子，遍地遊行，尋找可吞吃的人。

約翰二書 1:7-8
因為世上有許多迷惑人的出來，他們不認耶穌基督是成了肉身來的。這就是那迷惑人、敵基督的。 8 你們要\textcolor{red}{小心}，不要失去你們所做的工，乃要得著滿足的賞賜。




















\subsection{警醒}
\subsubsection{警醒禱告/做工/看门/等主/記念圣经教导/在真道上立穩/穿戴救恩/不被吞吃/堅固弱者/看守衣服}
\begin{table}[H]
\centering
\begin{tabular}{l|l|l|l}\hline
\multicolumn{2}{c|}{四福音} & \multicolumn{2}{c}{其他} \\\hline\hline
主哪一天來到& 太 24:32-44&警醒記念保罗三年晝夜不住流淚的勸誡& 徒 20:25-31\\\hline
迎接新郎,進去坐席& 太 25:1-13&警醒在真道上立穩，做大丈夫，要剛強&林前16:13 \\\hline
警醒禱告，免得入了迷惑& 太 26:37-41&恆切禱告警醒、感恩&西 4:2-4 \\\hline
警醒禱告/做工/看门，迎接主來& 可 13:33-37&不要睡覺要警醒謹守,穿戴信/愛/得救的盼望&贴前 5:6 \\\hline
警醒禱告，免得入了迷惑& 可 14:32-38&謹慎自守警醒禱告& 彼前4:7\\\hline
束腰點燈,等候主人從婚筵回來& 路 12:35-40&警醒被仇敵吞吃&彼前5:8 \\\hline
常常祈求逃避要来的,站立在人子前& 路 21:34-36&堅固那剩下將要衰微的，& 啟 3:2-3\\\hline
& &警醒看守衣服免得赤身而行&啟 16:15 \\\hline\hline
\end{tabular}
\end{table}

\subsubsection{详细经文}
馬太福音 24:32-44\\
你們可以從無花果樹學個比方。當樹枝發嫩長葉的時候，你們就知道夏天近了。 33 這樣，你們看見這一切的事，也該知道人子近了，正在門口了。 34 我實在告訴你們：這世代還沒有過去，這些事都要成就。 35 天地要廢去，我的話卻不能廢去。 36 但那日子、那時辰，沒有人知道，連天上的使者也不知道，子也不知道，唯獨父知道。 37 挪亞的日子怎樣，人子降臨也要怎樣。 38 當洪水以前的日子，人照常吃喝嫁娶，直到挪亞進方舟的那日， 39 不知不覺洪水來了，把他們全都沖去。人子降臨也要這樣。 40 那時，兩個人在田裡，取去一個，撇下一個； 41 兩個女人推磨，取去一個，撇下一個。 42 所以，你們要\textcolor{red}{警醒}，因為不知道你們的主是哪一天來到。 43 家主若知道幾更天有賊來，就必\textcolor{red}{警醒}，不容人挖透房屋，這是你們所知道的。 44 所以，你們也要預備，因為你們想不到的時候，人子就來了。

馬太福音 25:1-13\\
那時，天國好比十個童女拿著燈出去迎接新郎。 2 其中有五個是愚拙的，五個是聰明的。 3 愚拙的拿著燈，卻預備油； 4 聰明的拿著燈，又預備油在器皿裡。 5 新郎遲延的時候，她們都打盹睡著了。 6 半夜有人喊著說：『新郎來了，你們出來迎接他！』 7 那些童女就都起來收拾燈。 8 愚拙的對聰明的說：『請分點油給我們，因為我們的燈要滅了。』 9 聰明的回答說：『恐怕不夠你我用的，不如你們自己到賣油的那裡去買吧。』 10 她們去買的時候，新郎到了，那預備好了的同他進去坐席，門就關了。 11 其餘的童女隨後也來了，說：『主啊，主啊，給我們開門！』 12 他卻回答說：『我實在告訴你們：我不認識你們。』 13 所以，你們要\textcolor{red}{警醒}，因為那日子、那時辰，你們不知道。

馬太福音 26:36-41\\
你們坐在這裡，等我到那邊去禱告。」 37 於是帶著彼得和西庇太的兩個兒子同去，就憂愁起來，極其難過， 38 便對他們說：「我心裡甚是憂傷，幾乎要死。你們在這裡等候，和我一同\textcolor{red}{警醒}。」 39 他就稍往前走，俯伏在地，禱告說：「我父啊！倘若可行，求你叫這杯離開我！然而，不要照我的意思，只要照你的意思。」 40 來到門徒那裡，見他們睡著了，就對彼得說：「怎麼樣，你們不能同我\textcolor{red}{警醒}片時嗎？ 41 總要\textcolor{red}{警醒}禱告，免得入了迷惑。你們心靈固然願意，肉體卻軟弱了。」


马可福音 13:33-37\\
你们要谨慎，\textcolor{red}{警醒}祈祷，因为你们不晓得那日期几时来到。 34 这事正如一个人离开本家，寄居外邦，把权柄交给仆人，分派各人当做的工，又吩咐看门的\textcolor{red}{警醒}。 35 所以，你们要\textcolor{red}{警醒}，因为你们不知道家主什么时候来，或晚上，或半夜，或鸡叫，或早晨。 36 恐怕他忽然来到，看见你们睡着了。 37 我对你们所说的话，也是对众人说：要\textcolor{red}{警醒}！

马可福音 14:32-38\\
耶穌對門徒說：「你們坐在這裡，等我禱告。」 33 於是帶著彼得、雅各、約翰同去，就驚恐起來，極其難過， 34 對他們說：「我心裡甚是憂傷，幾乎要死。你們在這裡等候，\textcolor{red}{警醒}。」 35 他就稍往前走，俯伏在地，禱告說倘若可行，便叫那時候過去。 36 他說：「阿爸，父啊！在你凡事都能，求你將這杯撤去！然而，不要從我的意思，只要從你的意思。」 37 耶穌回來，見他們睡著了，就對彼得說：「西門，你睡覺嗎？不能\textcolor{red}{警醒}片時嗎？ 38 總要\textcolor{red}{警醒}禱告，免得入了迷惑。你們心靈固然願意，肉體卻軟弱了。」

路加福音 12:35-40\\
35 「你們腰裡要束上帶，燈也要點著， 36 自己好像僕人等候主人從婚姻的筵席上回來。他來到叩門，就立刻給他開門。 37 主人來了，看見僕人\textcolor{red}{警醒}，那僕人就有福了。我實在告訴你們：主人必叫他們坐席，自己束上帶進前伺候他們。 38 或是二更天來，或是三更天來，看見僕人這樣，那僕人就有福了。 39 家主若知道賊什麼時候來，就必\textcolor{red}{警醒}，不容賊挖透房屋，這是你們所知道的。 40 你們也要預備，因為你們想不到的時候，人子就來了。

路加福音 21:34-36\\
你们要谨慎，恐怕因贪食、醉酒并今生的思虑累住你们的心，那日子就如同网罗忽然临到你们， 35 因为那日子要这样临到全地上一切居住的人。 36 你们要时时\textcolor{red}{警醒}，常常祈求，使你们能逃避这一切要来的事，得以站立在人子面前。

使徒行傳 20:25-31\\
我素常在你們中間來往，傳講神國的道，如今我曉得，你們以後都不得再見我的面了。 26 所以我今日向你們證明，你們中間無論何人死亡，罪不在我身上， 27 因為神的旨意我並沒有一樣避諱不傳給你們的。 28 聖靈立你們做全群的監督，你們就當為自己謹慎，也為全群謹慎，牧養神的教會，就是他用自己血所買來的 d。 29 我知道我去之後，必有凶暴的豺狼進入你們中間，不愛惜羊群。 30 就是你們中間也必有人起來，說悖謬的話，要引誘門徒跟從他們。 31 所以你們應當\textcolor{red}{警醒}，記念我三年之久晝夜不住地流淚勸誡你們各人。

哥林多前書 16:13\\
你們務要\textcolor{red}{警醒}，在真道上站立得穩，要做大丈夫，要剛強。

歌羅西書 4:2-4
你們要恆切禱告，在此\textcolor{red}{警醒}、感恩。 3 也要為我們禱告，求神給我們開傳道的門，能以講基督的奧祕——我為此被捆鎖—— 4 叫我按著所該說的話將這奧祕發明出來。

帖撒羅尼迦前書 5:6-11\\
所以，我們不要睡覺，像別人一樣，總要\textcolor{red}{警醒}謹守。 7 因為睡了的人是在夜間睡，醉了的人是在夜間醉。 8 但我們既然屬乎白晝，就應當謹守，把信和愛當做護心鏡遮胸，把得救的盼望當做頭盔戴上。 9 因為神不是預定我們受刑，乃是預定我們藉著我們主耶穌基督得救。 10 他替我們死，叫我們無論\textcolor{red}{醒著}、睡著，都與他同活。 11 所以，你們該彼此勸慰，互相建立，正如你們素常所行的。

彼得前書 4:7\\
萬物的結局近了，所以你們要謹慎自守，\textcolor{red}{警醒}禱告。

警醒仇敵吞吃 - 彼得前書 5:8\\
務要謹守、\textcolor{red}{警醒}，因為你們的仇敵魔鬼如同吼叫的獅子，遍地遊行，尋找可吞吃的人。

啟示錄 3:2,3\\
你要\textcolor{red}{警醒}，堅固那剩下將要衰微的，因我見你的行為在我神面前，沒有一樣是完全的。…

警醒來看守衣服 - 啟示錄 16:15\\
看哪，我來像賊一樣。那\textcolor{red}{警醒}、看守衣服，免得赤身而行，叫人見他羞恥的有福了！








\subsection{利劍(兩刃)——別迦摩}

\subsubsection{以賽亞書 11:1-5, 49:1-13}
\begin{itemize} 
\itemsep-.25em
\item 從耶西的本必發一條，從他根生的枝子必結果實。 2 耶和華的靈必住在他身上，就是使他有智慧和聰明的靈，謀略和能力的靈，知識和敬畏耶和華的靈。 3 他必以敬畏耶和華為樂，行審判不憑眼見，斷是非也不憑耳聞； 4 卻要以公義審判貧窮人，以正直判斷世上的謙卑人，以口中的杖擊打世界，以嘴裡的氣殺戮惡人。 5 公義必當他的腰帶，信實必當他脅下的帶子。
\item 1 「眾海島啊，當聽我言！遠方的眾民哪，留心而聽！自我出胎，耶和華就選召我；自出母腹，他就提我的名。 2 他使我的口如快刀，將我藏在他手蔭之下；又使我成為磨亮的箭，將我藏在他箭袋之中。 3 對我說：『你是我的僕人以色列，我必因你得榮耀。』 4 我卻說：『我勞碌是徒然，我盡力是虛無虛空，然而我當得的理必在耶和華那裡，我的賞賜必在我神那裡。』5 「耶和華從我出胎造就我做他的僕人，要使雅各歸向他，使以色列到他那裡聚集。原來耶和華看我為尊貴，我的神也成為我的力量。 6 現在他說：『你做我的僕人，使雅各眾支派復興，使以色列中得保全的歸回，尚為小事，我還要使你做外邦人的光，叫你施行我的救恩，直到地極。』」 7 救贖主以色列的聖者耶和華，對那被人所藐視、本國所憎惡、官長所虐待的如此說：「君王要看見就站起，首領也要下拜，都因信實的耶和華，就是揀選你以色列的聖者。」8 耶和華如此說：「在悅納的時候，我應允了你；在拯救的日子，我濟助了你。我要保護你，使你做眾民的中保，復興遍地，使人承受荒涼之地為業。 9 對那被捆綁的人說：『出來吧！』對那在黑暗的人說：『顯露吧！』他們在路上必得飲食，在一切淨光的高處必有食物。 10 不飢不渴，炎熱和烈日必不傷害他們，因為憐恤他們的必引導他們，領他們到水泉旁邊。 11 我必使我的眾山成為大道，我的大路也被修高。 12 看哪，這些從遠方來，這些從北方，從西方來，這些從秦國來。」 13 諸天哪，應當歡呼！大地啊，應當快樂！眾山哪，應當發聲歌唱！因為耶和華已經安慰他的百姓，也要憐恤他困苦之民。
\end{itemize}

\subsubsection{以弗所書 6:17}
並戴上救恩的頭盔，拿著聖靈的寶劍，就是神的道。

\subsubsection{希伯來書 4:12-13}
12 神的道是活潑的，是有功效的，比一切兩刃的劍更快，甚至魂與靈、骨節與骨髓，都能刺入、剖開，連心中的思念和主意都能辨明； 13 並且被造的沒有一樣在他面前不顯然的，原來萬物在那與我們有關係的主眼前，都是赤露敞開的

\subsubsection{啟示錄 1:12-16, 2:12-17, 19:11-16}
\begin{itemize} 
\itemsep-.25em
\item 12 我轉過身來，要看是誰發聲與我說話；既轉過來，就看見七個金燈臺， 13 燈臺中間有一位好像人子，身穿長衣，直垂到腳，胸間束著金帶。 14 他的頭與髮皆白，如白羊毛，如雪，眼目如同火焰， 15 腳好像在爐中鍛煉光明的銅，聲音如同眾水的聲音。 16 他右手拿著七星，從他口中出來一把兩刃的利劍，面貌如同烈日放光。
\item 「你要寫信給別迦摩教會的使者說：『那有兩刃利劍的說： 13 我知道你的居所，就是有撒旦座位之處。當我忠心的見證人安提帕在你們中間、撒旦所住的地方被殺之時，你還堅守我的名，沒有棄絕我的道。 14 然而，有幾件事我要責備你，因為在你那裡，有人服從了巴蘭的教訓。這巴蘭曾教導巴勒將絆腳石放在以色列人面前，叫他們吃祭偶像之物，行姦淫的事。 15 你那裡也有人照樣服從了尼哥拉一黨人的教訓。『所以，你當悔改！若不悔改，我就快臨到你那裡，用我口中的劍攻擊他們。 17 聖靈向眾教會所說的話，凡有耳的，就應當聽！得勝的，我必將那隱藏的嗎哪賜給他，並賜他一塊白石，石上寫著新名，除了那領受的以外，沒有人能認識。』
\item 我觀看，見天開了。有一匹白馬，騎在馬上的稱為「誠信真實」，他審判、爭戰都按著公義。 12 他的眼睛如火焰，他頭上戴著許多冠冕，又有寫著的名字，除了他自己沒有人知道。 13 他穿著濺了血的衣服，他的名稱為「神之道」。 14 在天上的眾軍騎著白馬，穿著細麻衣，又白又潔，跟隨他。 15 有利劍從他口中出來，可以擊殺列國。他必用鐵杖轄管 b他們，並要踹全能神烈怒的酒榨。 16 在他衣服和大腿上有名寫著說：「萬王之王，萬主之主。」
\end{itemize} 





\subsection{龍}
\subsubsection{大紅龍要吞婦人的孩子}

天上又現出異象來：有一條大紅龍，七頭十角，七頭上戴著七個冠冕， 4 牠的尾巴拖拉著天上星辰的三分之一，摔在地上。龍就站在那將要生產的婦人面前，等她生產之後，要吞吃她的孩子。

\subsubsection{天上有爭戰龍就被摔於地}

7 在天上就有了爭戰，米迦勒同他的使者與龍爭戰。龍也同牠的使者去爭戰， 8 並沒有得勝，天上再沒有牠們的地方。 9 大龍就是那古蛇，名叫魔鬼，又叫撒旦，是迷惑普天下的。牠被摔在地上，牠的使者也一同被摔下去。 10 我聽見在天上有大聲音說：「我神的救恩、能力、國度並他基督的權柄，現在都來到了！因為那在我們神面前晝夜控告我們弟兄的，已經被摔下去了。 11 弟兄勝過牠，是因羔羊的血和自己所見證的道，他們雖至於死也不愛惜性命。 12 所以諸天和住在其中的，你們都快樂吧！只是地與海有禍了！因為魔鬼知道自己的時候不多，就氣憤憤地下到你們那裡去了。」










\subsection{嗎哪(隱藏的)——別迦摩}

\subsubsection{出埃及記 16:4-36}
耶和華對摩西說：「我要將糧食從天降給你們。百姓可以出去，每天收每天的份，我好試驗他們遵不遵我的法度。 5 到第六天，他們要把所收進來的預備好了，比每天所收的多一倍。」 6 摩西、亞倫對以色列眾人說：「到了晚上，你們要知道是耶和華將你們從埃及地領出來的。 7 早晨，你們要看見耶和華的榮耀，因為耶和華聽見你們向他所發的怨言了。我們算什麼，你們竟向我們發怨言呢？」 8 摩西又說：「耶和華晚上必給你們肉吃，早晨必給你們食物得飽，因為你們向耶和華發的怨言，他都聽見了。我們算什麼？你們的怨言不是向我們發的，乃是向耶和華發的。」 9 摩西對亞倫說：「你告訴以色列全會眾說：『你們就近耶和華面前，因為他已經聽見你們的怨言了。』」 10 亞倫正對以色列全會眾說話的時候，他們向曠野觀看，不料，耶和華的榮光在雲中顯現。 11 耶和華曉諭摩西說： 12 「我已經聽見以色列人的怨言。你告訴他們說：『到黃昏的時候，你們要吃肉，早晨必有食物得飽，你們就知道我是耶和華你們的神。』」13 到了晚上，有鵪鶉飛來，遮滿了營。早晨，在營四圍的地上有露水。 14 露水上升之後，不料，野地面上有如白霜的小圓物。 15 以色列人看見，不知道是什麼，就彼此對問說：「這是什麼呢？」摩西對他們說：「這就是耶和華給你們吃的食物。 16 耶和華所吩咐的是這樣：你們要按著各人的飯量，為帳篷裡的人按著人數收起來，各拿一俄梅珥。」 17 以色列人就這樣行，有多收的，有少收的。 18 及至用俄梅珥量一量，多收的也沒有餘，少收的也沒有缺，各人按著自己的飯量收取。 19 摩西對他們說：「所收的，不許什麼人留到早晨。」 20 然而他們不聽摩西的話，內中有留到早晨的，就生蟲變臭了。摩西便向他們發怒。21 他們每日早晨，按著各人的飯量收取，日頭一發熱，就消化了。 22 到第六天，他們收了雙倍的食物，每人兩俄梅珥。會眾的官長來告訴摩西， 23 摩西對他們說：「耶和華這樣說：『明天是聖安息日，是向耶和華守的聖安息日。你們要烤的就烤了，要煮的就煮了，所剩下的都留到早晨。』」 24 他們就照摩西的吩咐留到早晨，也不臭，裡頭也沒有蟲子。 25 摩西說：「你們今天吃這個吧，因為今天是向耶和華守的安息日，你們在田野必找不著了。 26 六天可以收取，第七天乃是安息日，那一天必沒有了。」 27 第七天，百姓中有人出去收，什麼也找不著。 28 耶和華對摩西說：「你們不肯守我的誡命和律法，要到幾時呢？ 29 你們看，耶和華既將安息日賜給你們，所以第六天他賜給你們兩天的食物，第七天各人要住在自己的地方，不許什麼人出去。」 30 於是百姓第七天安息了。31 這食物，以色列家叫嗎哪，樣子像芫荽子，顏色是白的，滋味如同攙蜜的薄餅。 32 摩西說：「耶和華所吩咐的是這樣：要將一滿俄梅珥嗎哪留到世世代代，使後人可以看見我當日將你們領出埃及地，在曠野所給你們吃的食物。」 33 摩西對亞倫說：「你拿一個罐子，盛一滿俄梅珥嗎哪，存在耶和華面前，要留到世世代代。」 34 耶和華怎麼吩咐摩西，亞倫就怎麼行，把嗎哪放在法櫃前存留。 35 以色列人吃嗎哪共四十年，直到進了有人居住之地，就是迦南的境界。 36 （俄梅珥乃伊法十分之一。）

\subsubsection{民數記 11:4-9}
4 他們中間的閒雜人大起貪慾的心。以色列人又哭號說：「誰給我們肉吃呢？ 5 我們記得在埃及的時候不花錢就吃魚，也記得有黃瓜、西瓜、韭菜、蔥、蒜。 6 現在我們的心血枯竭了，除這嗎哪以外，在我們眼前並沒有別的東西！」 7 這嗎哪彷彿芫荽子，又好像珍珠。 8 百姓周圍行走，把嗎哪收起來，或用磨推，或用臼搗，煮在鍋中，又做成餅，滋味好像新油。 9 夜間露水降在營中，嗎哪也隨著降下。


\subsubsection{申命記 8:3}
他苦煉你，任你飢餓，將你和你列祖所不認識的嗎哪賜給你吃，使你知道人活著不是單靠食物，乃是靠耶和華口裡所出的一切話。

\subsubsection{約書亞記 5:12}
他們吃了那地的出產，第二日嗎哪就止住了，以色列人也不再有嗎哪了。那一年，他們卻吃迦南地的出產。

\subsubsection{尼希米記 9:15,20}
15從天上賜下糧食充他們的飢，從磐石使水流出解他們的渴，又吩咐他們進去得你起誓應許賜給他們的地。20 你也賜下你良善的靈教訓他們，未嘗不賜嗎哪使他們糊口，並賜水解他們的渴。

\subsubsection{詩篇 78:18-25}
他們心中試探神，隨自己所欲的求食物，
19 並且妄論神說：「神在曠野豈能擺設筵席嗎？
20 他曾擊打磐石，使水湧出成了江河，他還能賜糧食嗎？還能為他的百姓預備肉嗎？」
21 所以耶和華聽見就發怒，有烈火向雅各燒起，有怒氣向以色列上騰。
22 因為他們不信服神，不倚賴他的救恩。
23 他卻吩咐天空，又敞開天上的門，
24 降嗎哪像雨給他們吃，將天上的糧食賜給他們。
25 各人吃大能者的食物，他賜下糧食，使他們飽足。

\subsubsection{約翰福音 6:30-33，47-58}
30 他們又說：「你行什麼神蹟，叫我們看見就信你？你到底做什麼事呢？ 31 我們的祖宗在曠野吃過嗎哪，如經上寫著說：『他從天上賜下糧來給他們吃。』」 32 耶穌說：「我實實在在地告訴你們：那從天上來的糧不是摩西賜給你們的，乃是我父將天上來的真糧賜給你們。 33 因為神的糧就是那從天上降下來賜生命給世界的。」 47 我實實在在地告訴你們：信的人有永生。 48 我就是生命的糧。 49 你們的祖宗在曠野吃過嗎哪，還是死了； 50 這是從天上降下來的糧，叫人吃了就不死。 51 我是從天上降下來生命的糧，人若吃這糧，就必永遠活著。我所要賜的糧就是我的肉，為世人之生命所賜的。」52 因此，猶太人彼此爭論說：「這個人怎能把他的肉給我們吃呢？」 53 耶穌說：「我實實在在地告訴你們：你們若不吃人子的肉，不喝人子的血，就沒有生命在你們裡面。 54 吃我肉、喝我血的人就有永生，在末日我要叫他復活。 55 我的肉真是可吃的，我的血真是可喝的。 56 吃我肉、喝我血的人常在我裡面，我也常在他裡面。 57 永活的父怎樣差我來，我又因父活著，照樣，吃我肉的人也要因我活著。 58 這就是從天上降下來的糧。吃這糧的人就永遠活著，不像你們的祖宗吃過嗎哪還是死了。」

\subsubsection{哥林多前書 10:3，4}
並且都吃了一樣的靈食， 4也都喝了一樣的靈水；所喝的是出於隨著他們的靈磐石，那磐石就是基督。

\subsubsection{希伯來書 9:4}
有金香爐，有包金的約櫃，櫃裡有盛嗎哪的金罐和亞倫發過芽的杖並兩塊約版，

\subsubsection{啟示錄 2:17}
聖靈向眾教會所說的話，凡有耳的，就應當聽！得勝的，我必將那隱藏的嗎哪賜給他，並賜他一塊白石，石上寫著新名，除了那領受的以外，沒有人能認識。





\subsection{末期/末後}
\subsubsection{但以理書 12：5-13}
5 我但以理觀看，見另有兩個人站立，一個在河這邊，一個在河那邊。 6 有一個問那站在河水以上穿細麻衣的說：「這奇異的事到幾時才應驗呢？」 7 我聽見那站在河水以上穿細麻衣的，向天舉起左右手，指著活到永遠的主起誓說：「要到一載，二載，半載，打破聖民權力的時候，這一切事就都應驗了。」 8 我聽見這話，卻不明白，就說：「我主啊，這些事的\underline{結局אַחֲרִ֖ית}是怎樣呢？」 9 他說：「但以理啊，你只管去，因為這話已經隱藏封閉，直到\underline{קֵץ末時}。 10 必有許多人使自己清淨潔白，且被熬煉，但惡人仍必行惡。一切惡人都不明白，唯獨智慧人能明白。 11 從除掉常獻的燔祭，並設立那行毀壞可憎之物的時候，必有一千二百九十日。 12 等到一千三百三十五日的，那人便為有福！ 13 你且去等候\underline{結局}，因為你必安歇。到了\underline{末期קֵץ}，你必起來，享受你的福分。」

\subsubsection{馬太福音 10：5-42}
5 耶穌差這十二個人去，吩咐他們說：「外邦人的路你們不要走，撒馬利亞人的城你們不要進， 6 寧可往以色列家迷失的羊那裡去。 7 隨走隨傳，說『天國近了』。 8 醫治病人，叫死人復活，叫長大痲瘋的潔淨，把鬼趕出去。你們白白地得來，也要白白地捨去。 9 腰袋裡不要帶金銀銅錢， 10 行路不要帶口袋，不要帶兩件褂子，也不要帶鞋和拐杖，因為工人得飲食是應當的。 11 你們無論進哪一城、哪一村，要打聽那裡誰是好人，就住在他家，直住到走的時候。 12 進他家裡去，要請他的安。 13 那家若配得平安，你們所求的平安就必臨到那家；若不配得，你們所求的平安仍歸你們。 14 凡不接待你們、不聽你們話的人，你們離開那家或是那城的時候，就把腳上的塵土跺下去。 15 我實在告訴你們：當審判的日子，所多瑪和蛾摩拉所受的比那城還容易受呢！21 「弟兄要把弟兄，父親要把兒子，送到死地；兒女要與父母為敵，害死他們。 22 並且你們要為我的名被眾人恨惡，唯有忍耐到\underline{底τέλος}的，必然得救。 23 有人在這城裡逼迫你們，就逃到那城裡去。我實在告訴你們：以色列的城邑你們還沒有走遍，人子就到了。24 「學生不能高過先生，僕人不能高過主人。 25 學生和先生一樣，僕人和主人一樣，也就罷了。人既罵家主是別西卜 a，何況他的家人呢？ 26 所以不要怕他們，因為掩蓋的事沒有不露出來的，隱藏的事沒有不被人知道的。 27 我在暗中告訴你們的，你們要在明處說出來；你們耳中所聽的，要在房上宣揚出來。 28 那殺身體不能殺靈魂的，不要怕他們；唯有能把身體和靈魂都滅在地獄裡的，正要怕他。 29 兩個麻雀不是賣一分銀子嗎？若是你們的父不許，一個也不能掉在地上。 30 就是你們的頭髮，也都被數過了。 31 所以，不要懼怕，你們比許多麻雀還貴重。 32 凡在人面前認我的，我在我天上的父面前也必認他； 33 凡在人面前不認我的，我在我天上的父面前也必不認他。34 「你們不要想我來是叫地上太平，我來並不是叫地上太平，乃是叫地上動刀兵。 35 因為我來是叫『人與父親生疏，女兒與母親生疏，媳婦與婆婆生疏； 36 人的仇敵就是自己家裡的人』。 37 愛父母過於愛我的，不配做我的門徒；愛兒女過於愛我的，不配做我的門徒； 38 不背著他的十字架跟從我的，也不配做我的門徒。 39 得著生命的，將要失喪生命；為我失喪生命的，將要得著生命.40 「人接待你們，就是接待我；接待我，就是接待那差我來的。 41 人因為先知的名接待先知，必得先知所得的賞賜；人因為義人的名接待義人，必得義人所得的賞賜。 42 無論何人，因為門徒的名，只把一杯涼水給這小子裡的一個喝，我實在告訴你們：這人不能不得賞賜。」

\subsubsection{馬太福音 24：4-14}
 4 耶穌回答說：「你們要謹慎，免得有人迷惑你們。 5 因為將來有好些人冒我的名來，說：『我是基督』，並且要迷惑許多人。 6 你們也要聽見打仗和打仗的風聲，總不要驚慌，因為這些事是必須有的，只是\underline{末期 τέλος}還沒有到。 7 民要攻打民，國要攻打國，多處必有饑荒、地震， 8 這都是災難的起頭。 9 那時，人要把你們陷在患難裡，也要殺害你們；你們又要為我的名被萬民恨惡。10 「那時，必有許多人跌倒，也要彼此陷害、彼此恨惡， 11 且有好些假先知起來，迷惑多人。 12 只因不法的事增多，許多人的愛心才漸漸冷淡了。 13 唯有忍耐到\underline{底τέλος}的，必然得救。 14 這天國的福音要傳遍天下，對萬民作見證，然後\underline{末期 τέλος}才來到。

\subsubsection{馬可福音 13：3-13}
3 耶穌在橄欖山上對聖殿而坐，彼得、雅各、約翰和安得烈暗暗地問他說： 4 「請告訴我們，什麼時候有這些事呢？這一切事將成的時候有什麼預兆呢？」 5 耶穌說：「你們要謹慎，免得有人迷惑你們。 6 將來有好些人冒我的名來，說『我是基督』，並且要迷惑許多人。 7 你們聽見打仗和打仗的風聲，不要驚慌，這些事是必須有的，只是\underline{末期τέλος}還沒有到。 8 民要攻打民，國要攻打國，多處必有地震、饑荒，這都是災難的起頭。9 「但你們要謹慎，因為人要把你們交給公會，並且你們在會堂裡要受鞭打，又為我的緣故站在諸侯與君王面前，對他們作見證。 10 然而，福音必須先傳給萬民。 11 人把你們拉去交官的時候，不要預先思慮說什麼，到那時候，賜給你們什麼話，你們就說什麼；因為說話的不是你們，乃是聖靈。 12 弟兄要把弟兄，父親要把兒子，送到死地；兒女要起來與父母為敵，害死他們。 13 並且你們要為我的名被眾人恨惡，唯有忍耐到\underline{底τέλος}的，必然得救。

\subsubsection{路加福音 21：6-9}
7 他們問他說：「夫子，什麼時候有這事呢？這事將到的時候有什麼預兆呢？」 8 耶穌說：「你們要謹慎，不要受迷惑。因為將來有好些人冒我的名來，說『我是基督』，又說『時候近了』。你們不要跟從他們！ 9 你們聽見打仗和擾亂的事，不要驚惶，因為這些事必須先有，只是\underline{末期τέλος}不能立時就到。」
\subsubsection{路加福音 21：20-28}
20 「你們看見耶路撒冷被兵圍困，就可知道它成荒場的日子近了。 21 那時，在猶太的，應當逃到山上；在城裡的，應當出來；在鄉下的，不要進城。 22 因為這是報應的日子，使經上所寫的都得應驗。 23 當那些日子，懷孕的和奶孩子的有禍了！因為將有大災難降在這地方，也有震怒臨到這百姓。 24 他們要倒在刀下，又被擄到各國去。耶路撒冷要被外邦人踐踏，直到外邦人的日期滿了。 25 日、月、星辰要顯出異兆，地上的邦國也有困苦，因海中波浪的響聲就慌慌不定。 26 天勢都要震動，人想起那將要臨到世界的事，就都嚇得魂不附體。 27 那時，他們要看見人子有能力、有大榮耀駕雲降臨。 28 一有這些事，你們就當挺身昂首，因為你們\underline{得贖的日子}近了。」

\subsubsection{哥林多前書 1：4-9}
4 我常為你們感謝我的神，因神在基督耶穌裡所賜給你們的恩惠， 5 又因你們在他裡面凡事富足，口才、知識都全備。 6 正如我為基督作的見證在你們心裡得以堅固， 7 以致你們在恩賜上沒有一樣不及人的，等候我們的主耶穌基督顯現。 8 他也必堅固你們到\underline{底τέλους}，叫你們在我們主耶穌基督的日子無可責備。 9 神是信實的，你們原是被他所召，好與他兒子我們的主耶穌基督一同得份。

\subsubsection{哥林多前書 10：4-13}
1 弟兄們，我不願意你們不曉得：我們的祖宗從前都在雲下，都從海中經過， 2 都在雲裡、海裡受洗歸了摩西， 3 並且都吃了一樣的靈食， 4 也都喝了一樣的靈水；所喝的是出於隨著他們的靈磐石，那磐石就是基督。 5 但他們中間多半是神不喜歡的人，所以在曠野倒斃。 6 這些事都是我們的鑒戒，叫我們不要貪戀惡事，像他們那樣貪戀的。 7 也不要拜偶像，像他們有人拜的，如經上所記：「百姓坐下吃喝，起來玩耍。」 8 我們也不要行姦淫，像他們有人行的，一天就倒斃了二萬三千人。 9 也不要試探主，像他們有人試探的，就被蛇所滅。 10 你們也不要發怨言，像他們有發怨言的，就被滅命的所滅。 11 他們遭遇這些事都要作為鑒戒，並且寫在經上正是警戒我們這\underline{末世τέλη}的人。 12 所以，自己以為站得穩的須要謹慎，免得跌倒。 13 你們所遇見的試探，無非是人所能受的。神是信實的，必不叫你們受試探過於所能受的，在受試探的時候，總要給你們開一條出路，叫你們能忍受得住。

\subsubsection{哥林多前書 15：20-28}
20 但基督已經從死裡復活，成為睡了之人初熟的果子。 21 死既是因一人而來，死人復活也是因一人而來。 22 在亞當裡眾人都死了，照樣，在基督裡眾人也都要復活。 23 但各人是按著自己的次序復活：初熟的果子是基督，以後在他來的時候，是那些屬基督的。 24 再後\underline{末期τέλος}到了，那時基督既將一切執政的、掌權的、有能的都毀滅了，就把國交於父神。 25 因為基督必要做王，等神把一切仇敵都放在他的腳下。 26 儘末了所毀滅的仇敵就是死， 27 因為經上說：「神叫萬物都服在他的腳下。」既說萬物都服了他，明顯那叫萬物服他的不在其內了。 28 萬物既服了他，那時子也要自己服那叫萬物服他的，叫神在萬物之上、為萬物之主。

\subsubsection{希伯來書 6：1-8}
1 所以，我們應當離開基督道理的開端，竭力進到完全的地步，不必再立根基，就如那懊悔死行、信靠神、 2 各樣洗禮、按手之禮、死人復活，以及永遠審判各等教訓。 3 神若許我們，我們必如此行。 4 論到那些已經蒙了光照，嘗過天恩的滋味，又於聖靈有份， 5 並嘗過神善道的滋味，覺悟來世權能的人， 6 若是離棄道理，就不能叫他們重新懊悔了。因為他們把神的兒子重釘十字架，明明地羞辱他。 7 就如一塊田地，吃過屢次下的雨水，生長菜蔬，合乎耕種的人用，就從神得福； 8 若長荊棘和蒺藜，必被廢棄，近於咒詛，\underline{結局 τέλους}就是焚燒。

\subsubsection{彼得前書 4：7-19}
7 萬物的\underline{結局τέλος}近了，所以你們要謹慎自守，警醒禱告。 8 最要緊的是彼此切實相愛，因為愛能遮掩許多的罪。 9 你們要互相款待，不發怨言。 10 各人要照所得的恩賜彼此服侍，做神百般恩賜的好管家。 11 若有講道的，要按著神的聖言講；若有服侍人的，要按著神所賜的力量服侍，叫神在凡事上因耶穌基督得榮耀——原來榮耀、權能都是他的，直到永永遠遠！阿們。12 親愛的弟兄啊，有火煉的試驗臨到你們，不要以為奇怪，似乎是遭遇非常的事； 13 倒要歡喜，因為你們是與基督一同受苦，使你們在他榮耀顯現的時候，也可以歡喜快樂。 14 你們若為基督的名受辱罵，便是有福的，因為神榮耀的靈常住在你們身上。 15 你們中間卻不可有人因為殺人、偷竊、作惡、好管閒事而受苦。 16 若為做基督徒受苦，卻不要羞恥，倒要因這名歸榮耀給神。 17 因為時候到了，審判要從神的家起首。若是先從我們起首，那不信從神福音的人將有何等的\underline{結局τέλος}呢？ 18 「若是義人僅僅得救，那不虔敬和犯罪的人將有何地可站呢？」 19 所以，那照神旨意受苦的人要一心為善，將自己靈魂交於那信實的造化之主。

\subsubsection{啟示錄 2：24-29}

24 『至於你們推雅推喇其餘的人，就是一切不從那教訓、不曉得他們素常所說撒旦深奧之理的人，我告訴你們：我不將別的擔子放在你們身上； 25 但你們已經有的，總要持守，直等到\underline{（the end）τέλους}我來。 26 那得勝又遵守我命令到底的，我要賜給他權柄制伏列國， 27 他必用鐵杖轄管 b他們，將他們如同窯戶的瓦器打得粉碎，像我從我父領受的權柄一樣。 28 我又要把晨星賜給他。 29 聖靈向眾教會所說的話，凡有耳的，就應當聽！』

\subsubsection{啟示錄 21：6-8}
他又對我說：「都成了！我是阿拉法，我是俄梅戛；我是初，我是\underline{終 τέλος}。我要將生命泉的水白白賜給那口渴的人喝。 7 得勝的，必承受這些為業：我要做他的神，他要做我的兒子。 8 唯有膽怯的、不信的、可憎的、殺人的、淫亂的、行邪術的、拜偶像的和一切說謊話的，他們的份就在燒著硫磺的火湖裡，這是第二次的死。」

\subsubsection{啟示錄 22：12-15}
「看哪，我必快來！賞罰在我，要照各人所行的報應他。 13 我是阿拉法，我是俄梅戛；我是首先的，我是\underline{末後τέλος}的；我是初，我是終。 14 那些洗淨自己衣服的有福了！可得權柄能到生命樹那裡，也能從門進城。 15 城外有那些犬類、行邪術的、淫亂的、殺人的、拜偶像的，並一切喜好說謊言、編造虛謊的。











\subsection{那日}

\subsubsection{以賽亞書 2：2-22}
2\underline{ 末後的日子}，耶和華殿的山必堅立超乎諸山，高舉過於萬嶺，萬民都要流歸這山。 3 必有許多國的民前往，說：「來吧！我們登耶和華的山，奔雅各神的殿。主必將他的道教訓我們，我們也要行他的路，因為訓誨必出於錫安，耶和華的言語必出於耶路撒冷。」 4 他必在列國中施行審判，為許多國民斷定是非。他們要將刀打成犁頭，把槍打成鐮刀。這國不舉刀攻擊那國，他們也不再學習戰事。5 雅各家啊，來吧，我們在耶和華的光明中行走！ 6 耶和華，你離棄了你百姓雅各家，是因他們充滿了東方的風俗，做觀兆的，像非利士人一樣，並與外邦人擊掌。 7 他們的國滿了金銀，財寶也無窮；他們的地滿了馬匹，車輛也無數。 8 他們的地滿了偶像，他們跪拜自己手所造的，就是自己指頭所做的。 9 卑賤人屈膝，尊貴人下跪，所以不可饒恕他們。 10 你當進入巖穴，藏在土中，躲避耶和華的驚嚇和他威嚴的榮光。 11 到\underline{ 那日}，眼目高傲的必降為卑，性情狂傲的都必屈膝，唯獨耶和華被尊崇。12 必有萬軍耶和華\underline{ 降罰的一個日子}，要臨到驕傲狂妄的，一切自高的都必降為卑； 13 又臨到黎巴嫩高大的香柏樹和巴珊的橡樹； 14 又臨到一切高山的峻嶺； 15 又臨到高臺和堅固城牆； 16 又臨到他施的船隻，並一切可愛的美物。 17 驕傲的必屈膝，狂妄的必降卑，在\underline{ 那日}，唯獨耶和華被尊崇。18 偶像必全然廢棄。 19 耶和華興起使地大震動的時候，人就進入石洞，進入土穴，躲避耶和華的驚嚇和他威嚴的榮光。 20 到那日，人必將為拜而造的金偶像、銀偶像拋給田鼠和蝙蝠。 21 到耶和華興起使地大震動的時候，人好進入磐石洞中和巖石穴裡，躲避耶和華的驚嚇和他威嚴的榮光。 22 你們休要倚靠世人，他鼻孔裡不過有氣息，他在一切事上可算什麼呢？


\subsubsection{以賽亞書 3：16-25}
 耶和華又說：「因為錫安的女子狂傲，行走挺項，賣弄眼目，俏步徐行，腳下玎璫， 17 所以主必使錫安的女子頭長禿瘡，耶和華又使她們赤露下體。」 18 到 \underline{那日}，主必除掉她們華美的腳釧、髮網、月牙圈、 19 耳環、手鐲、蒙臉的帕子、 20 華冠、足鏈、華帶、香盒、符囊、 21 戒指、鼻環、 22 吉服、外套、雲肩、荷包、 23 手鏡、細麻衣、裹頭巾、蒙身的帕子。 24 必有臭爛代替馨香，繩子代替腰帶，光禿代替美髮，麻衣繫腰代替華服，烙傷代替美容。 25 你的男丁必倒在刀下，你的勇士必死在陣上。 26 錫安的城門必悲傷哀號，她必荒涼坐在地上。

\subsubsection{以賽亞書 4：1-6}
1 在 \underline{那日}，七個女人必拉住一個男人，說：「我們吃自己的食物，穿自己的衣服，但求你許我們歸你名下，求你除掉我們的羞恥！」2 到 \underline{那日}，耶和華發生的苗必華美尊榮，地的出產必為以色列逃脫的人顯為榮華茂盛。 3 主以公義的靈和焚燒的靈，將錫安女子的汙穢洗去，又將耶路撒冷中殺人的血除淨。那時，剩在錫安留在耶路撒冷的，就是一切住耶路撒冷在生命冊上記名的，必稱為聖。 4  5 耶和華也必在錫安全山並各會眾以上，使白日有煙雲，黑夜有火焰的光，因為在全榮耀之上必有遮蔽。 6 必有亭子，白日可以得蔭避暑，也可以作為藏身之處，躲避狂風暴雨
\subsubsection{以賽亞書 5：24-30}
火苗怎樣吞滅碎秸，乾草怎樣落在火焰之中，照樣，他們的根必像朽物，他們的花必像灰塵飛騰。因為他們厭棄萬軍之耶和華的訓誨，藐視以色列聖者的言語。 25 所以耶和華的怒氣向他的百姓發作，他的手伸出攻擊他們，山嶺就震動，他們的屍首在街市上好像糞土。雖然如此，他的怒氣還未轉消，他的手仍伸不縮。26 他必豎立大旗招遠方的國民，發嘶聲叫他們從地極而來。看哪，他們必急速奔來！ 27 其中沒有疲倦的、絆跌的，沒有打盹的、睡覺的，腰帶並不放鬆，鞋帶也不折斷。 28 他們的箭快利，弓也上了弦；馬蹄算如堅石，車輪好像旋風。 29 他們要吼叫像母獅子，咆哮像少壯獅子；他們要咆哮抓食，坦然叼去，無人救回。 30 \underline{那日}，他們要向以色列人吼叫，像海浪砰訇。人若望地，只見黑暗艱難，光明在雲中變為昏暗。

\subsubsection{以賽亞書 6：9-13}
撒拉弗說：「你去告訴這百姓說：『你們聽是要聽見，卻不明白；看是要看見，卻不曉得。』 10 要使這百姓心蒙脂油，耳朵發沉，眼睛昏迷，恐怕眼睛看見，耳朵聽見，心裡明白，回轉過來，便得醫治。」 11 我就說：「主啊，這到幾時為止呢？」他說：「\underline{直到}城邑荒涼，無人居住，房屋空閒無人，地土極其荒涼。 12 並且耶和華將人遷到遠方，在這境內撇下的地土很多。 13 境內剩下的人，若還有十分之一，也必被吞滅。像栗樹、橡樹，雖被砍伐，樹不子卻仍存留，這聖潔的種類在國中也是如此。」

\subsubsection{以賽亞書 7：17-25}
17 耶和華必使亞述王\underline{攻擊你的日子}臨到你和你的百姓並你的父家，自從以法蓮離開猶大以來，\underline{未曾有這樣的日子}。18 「\underline{那時}，耶和華要發嘶聲，使埃及江河源頭的蒼蠅和亞述地的蜂子飛來。 19 都必飛來，落在荒涼的谷內，磐石的穴裡，和一切荊棘籬笆中，並一切的草場上。20 「那時，主必用大河外賃的剃頭刀，就是亞述王，剃去頭髮和腳上的毛，並要剃淨鬍鬚。21 「\underline{那時}，一個人要養活一隻母牛犢，兩隻母綿羊。 22 因為出的奶多，他就得吃奶油。在境內所剩的人，都要吃奶油與蜂蜜。23 「從前凡種一千棵葡萄樹，值銀一千舍客勒的地方，到那時，必長荊棘和蒺藜。 24 人上那裡去，必帶弓箭，因為遍地滿了荊棘和蒺藜。 25 所有用鋤刨挖的山地，你因怕荊棘和蒺藜，不敢上那裡去，只可成了放牛之處，為羊踐踏之地。」

\subsubsection{以賽亞書 8：5-22}
5 耶和華又曉諭我說： 6 「這百姓既厭棄西羅亞緩流的水，喜悅利汛和利瑪利的兒子， 7 因此主必使大河翻騰的水猛然沖來，就是亞述王和他所有的威勢。必漫過一切的水道，漲過兩岸； 8 必沖入猶大，漲溢氾濫，直到頸項。以馬內利啊！他展開翅膀，遍滿你的地。」9 列國的人民哪，任憑你們喧嚷，終必破壞！遠方的眾人哪，當側耳而聽！任憑你們束起腰來，終必破壞，你們束起腰來，終必破壞！ 10 任憑你們同謀，終歸無有，任憑你們言定，終不成立，因為神與我們同在。 11 耶和華以大能的手，指教我不可行這百姓所行的道，對我這樣說： 12 「這百姓說同謀背叛，你們不要說同謀背叛。他們所怕的，你們不要怕，也不要畏懼。 13 但要尊萬軍之耶和華為聖，以他為你們所當怕的，所當畏懼的。 14 他必作為聖所，卻向以色列兩家做絆腳的石頭、跌人的磐石，向耶路撒冷的居民作為圈套和網羅。 15 許多人必在其上絆腳跌倒，而且跌碎，並陷入網羅被纏住。」16 你要捲起律法書，在我門徒中間封住訓誨。 17 我要等候那掩面不顧雅各家的耶和華，我也要仰望他。 18 看哪，我與耶和華所給我的兒女，就是從住在錫安山萬軍之耶和華來的，在以色列中作為預兆和奇蹟.19 有人對你們說：「當求問那些交鬼的和行巫術的，就是聲音綿蠻、言語微細的。」你們便回答說：「百姓不當求問自己的神嗎？豈可為活人求問死人呢？ 20 人當以訓誨和法度為標準。」他們所說的若不與此相符，必不得見晨光。 21 他們必經過這地，受艱難，受飢餓，飢餓的時候心中焦躁，咒罵自己的君王和自己的神。 22 仰觀上天，俯察下地，不料，盡是艱難黑暗和幽暗的痛苦，他們必被趕入烏黑的黑暗中去。


\subsubsection{以賽亞書 9：11-21}
11 因此，耶和華要高舉利汛的敵人來攻擊以色列，並要激動以色列的仇敵， 12 東有亞蘭人，西有非利士人，他們張口要吞吃以色列。雖然如此，耶和華的怒氣還未轉消，他的手仍伸不縮。13 這百姓還沒有歸向擊打他們的主，也沒有尋求萬軍之耶和華。 14 因此，耶和華\underline{一日之間}，必從以色列中剪除頭與尾，棕枝與蘆葦。 15 長老和尊貴人就是頭，以謊言教人的先知就是尾。 16 因為引導這百姓的使他們走錯了路，被引導的都必敗亡。 17 所以主必不喜悅他們的少年人，也不憐恤他們的孤兒寡婦，因為各人是褻瀆的，是行惡的，並且各人的口都說愚妄的話。雖然如此，耶和華的怒氣還未轉消，他的手仍伸不縮。邪惡像火焚燒，燒滅荊棘和蒺藜，在稠密的樹林中著起來，就成為煙柱，旋轉上騰。 19 因萬軍之耶和華的烈怒，地都燒遍，百姓成為火柴，無人憐愛弟兄。 20 有人右邊搶奪，仍受飢餓；左邊吞吃，仍不飽足。各人吃自己膀臂上的肉。 21 瑪拿西吞吃以法蓮，以法蓮吞吃瑪拿西，又一同攻擊猶大。雖然如此，耶和華的怒氣還未轉消，他的手仍伸不縮

\subsubsection{以賽亞書 10：12-34}
 12 主在錫安山和耶路撒冷成就他一切工作的時候，主說：「我必罰亞述王自大的心和他高傲眼目的榮耀。 13 因為他說：『我所成就的事，是靠我手的能力和我的智慧，我本有聰明。我挪移列國的地界，搶奪他們所積蓄的財寶，並且我像勇士，使坐寶座的降為卑。 14 我的手夠到列國的財寶，好像人夠到鳥窩；我也得了全地，好像人拾起所棄的雀蛋。沒有動翅膀的，沒有張嘴的，也沒有鳴叫的。』」15 斧豈可向用斧砍木的自誇呢？鋸豈可向用鋸的自大呢？好比棍掄起那舉棍的，好比杖舉起那非木的人。 16 因此，主萬軍之耶和華必使亞述王的肥壯人變為瘦弱，在他的榮華之下必有火著起，如同焚燒一樣。 17 以色列的光必如火，他的聖者必如火焰，在一日之間，將亞述王的荊棘和蒺藜焚燒淨盡。 18 又將他樹林和肥田的榮耀全然燒盡，好像拿軍旗的昏過去一樣。 19 他林中剩下的樹必稀少，就是孩子也能寫其數。 到\underline{那日}，以色列所剩下的和雅各家所逃脫的，不再倚靠那擊打他們的，卻要誠實倚靠耶和華以色列的聖者。 21 所剩下的，就是雅各家所剩下的，必歸回全能的神。 22 以色列啊，你的百姓雖多如海沙，唯有剩下的歸回。原來滅絕的事已定，必有公義施行，如水漲溢。 23 因為主萬軍之耶和華在全地之中，必成就所定規的結局。24 所以主萬軍之耶和華如此說：「住錫安我的百姓啊，亞述王雖然用棍擊打你，又照埃及的樣子舉杖攻擊你，你卻不要怕他。 25 因為還有一點點時候，向你們發的憤恨就要完畢，我的怒氣要向他發作，使他滅亡。」 26 萬軍之耶和華要興起鞭來攻擊他，好像在俄立磐石那裡殺戮米甸人一樣；耶和華的杖要向海伸出，把杖舉起，像在埃及一樣。 27 到\underline{那日}，亞述王的重擔必離開你的肩頭，他的軛必離開你的頸項，那軛也必因肥壯的緣故撐斷。亞述王來到亞葉，經過米磯崙，在密抹安放輜重。 29 他們過了隘口，在迦巴住宿。拉瑪人戰兢，掃羅的基比亞人逃跑。 30 迦琳的居民 c哪，要高聲呼喊！萊煞人哪，須聽！哀哉，困苦的亞拿突啊！ 31 瑪得米那人躲避，基柄的居民逃遁。 32 當\underline{那日}，亞述王要在挪伯歇兵，向錫安女子的山，就是耶路撒冷的山，掄手攻他。33 看哪，主萬軍之耶和華以驚嚇削去樹枝，長高的必被砍下，高大的必被伐倒。 34 稠密的樹林，他要用鐵器砍下，黎巴嫩的樹木必被大能者伐倒。

\subsubsection{以賽亞書 11:1-5，10-16}
1 從耶西的本 a必發一條，從他根生的枝子必結果實。 2 耶和華的靈必住在他身上，就是使他有智慧和聰明的靈，謀略和能力的靈，知識和敬畏耶和華的靈。 3 他必以敬畏耶和華為樂，行審判不憑眼見，斷是非也不憑耳聞； 4 卻要以公義審判貧窮人，以正直判斷世上的謙卑人，以口中的杖擊打世界，以嘴裡的氣殺戮惡人。 5 公義必當他的腰帶，信實必當他脅下的帶子。到\underline{那日}，耶西的根立做萬民的大旗，外邦人必尋求他，他安息之所大有榮耀11 當\underline{那日}，主必二次伸手救回自己百姓中所餘剩的，就是在亞述、埃及、巴忒羅、古實、以攔、示拿、哈馬並眾海島所剩下的。 12 他必向列國豎立大旗，招回以色列被趕散的人，又從地的四方聚集分散的猶大人。 13 以法蓮的嫉妒就必消散，擾害猶大的必被剪除；以法蓮必不嫉妒猶大，猶大也不擾害以法蓮。 14 他們要向西飛，撲在非利士人的肩頭上，一同擄掠東方人，伸手按住以東和摩押，亞捫人也必順服他們。 15 耶和華必使埃及海汊枯乾，掄手用暴熱的風使大河分為七條，令人過去不致濕腳。 16 為主餘剩的百姓，就是從亞述剩下回來的，必有一條大道，如當日以色列從埃及地上來一樣。

\subsubsection{以賽亞書 13:2-18}
2 「應當在淨光的山豎立大旗，向群眾揚聲招手，使他們進入貴胄的門！ 3 我吩咐我所挑出來的人，我招呼我的勇士，就是那矜誇高傲之輩，為要成就我怒中所定的。」 4 山間有多人的聲音，好像是大國人民。有許多國的民聚集鬨嚷的聲音，這是萬軍之耶和華點齊軍隊，預備打仗。 5 他們從遠方來，從天邊來，就是耶和華並他惱恨的兵器，要毀滅這全地。6 你們要哀號，因為\underline{耶和華的日子}臨近了！這日來到，好像毀滅從全能者來到。 7 所以人手都必軟弱，人心都必消化。 8 他們必驚惶悲痛，愁苦必將他們抓住，他們疼痛好像產難的婦人一樣。彼此驚奇相看，臉如火焰。 9 \underline{耶和華的日子}臨到，必有殘忍、憤恨、烈怒，使這地荒涼，從其中除滅罪人。 10 天上的眾星群宿都不發光，日頭一出就變黑暗，月亮也不放光。 11 「我必因邪惡刑罰世界，因罪孽刑罰惡人，使驕傲人的狂妄止息，制伏強暴人的狂傲。 12 我必使人比精金還少，使人比俄斐純金更少。 13 我萬軍之耶和華在憤恨中發烈怒的日子，必使天震動，使地搖撼離其本位。 14 人必像被追趕的鹿，像無人收聚的羊，各歸回本族，各逃到本土。 15 凡被仇敵追上的必被刺死，凡被捉住的必被刀殺。 16 他們的嬰孩必在他們眼前摔碎，他們的房屋必被搶奪，他們的妻子必被玷汙。17 「我必激動瑪代人來攻擊他們，瑪代人不注重銀子，也不喜愛金子。 18 他們必用弓擊碎少年人，不憐憫婦人所生的，眼也不顧惜孩子。

\subsubsection{以賽亞書 17:4-9，12-14}
4 「到\underline{那日}，雅各的榮耀必致枵薄，他肥胖的身體必漸瘦弱。 5 就必像收割的人收斂禾稼，用手割取穗子，又像人在利乏音谷拾取遺落的穗子。 6 其間所剩下的不多，好像人打橄欖樹，在儘上的枝梢上只剩兩三個果子，在多果樹的旁枝上只剩四五個果子。」這是耶和華以色列的神說的。 7 當\underline{那日}，人必仰望造他們的主，眼目重看以色列的聖者。 8 他們必不仰望祭壇，就是自己手所築的，也不重看自己指頭所做的，無論是木偶是日像。 9 在\underline{那日}，他們的堅固城必像樹林中和山頂上所撇棄的地方，就是從前在以色列人面前被人撇棄的，這樣地就荒涼了....唉！多民鬨嚷，好像海浪砰訇；列邦奔騰，好像猛水滔滔。 13 列邦奔騰，好像多水滔滔，但神斥責他們，他們就遠遠逃避，又被追趕，如同山上的風前糠，又如暴風前的旋風土。 14 到\underline{晚上}有驚嚇，\underline{未到早晨}他們就沒有了。這是擄掠我們之人所得的份，是搶奪我們之人的報應。

\subsubsection{耶利米書 51:1-24}
耶和華如此說：「我必使毀滅的風颳起，攻擊巴比倫和住在立加米的人。 2 我要打發外邦人來到巴比倫，簸揚她，使她的地空虛，在她、\underline{遭禍的日子}他們要周圍攻擊她。 3 拉弓的要向拉弓的和貫甲挺身的射箭，不要憐惜她的少年人，要滅盡她的全軍。 4 他們必在迦勒底人之地被殺仆倒，在巴比倫的街上被刺透。」5 以色列和猶大雖然境內充滿違背以色列聖者的罪，卻沒有被他的神萬軍之耶和華丟棄。 6 你們要從巴比倫中逃奔，各救自己的性命！不要陷在她的罪孽中一同滅亡，因為這是耶和華報仇的時候，他必向巴比倫施行報應。 7 巴比倫素來是耶和華手中的金杯，使天下沉醉，萬國喝了她的酒就顛狂了。 8 巴比倫忽然傾覆毀壞，要為她哀號，為止她的疼痛，拿乳香或者可以治好。 9 我們想醫治巴比倫，她卻沒有治好。離開她吧！我們各人歸回本國，因為她受的審判通於上天，達到穹蒼。 10 耶和華已經彰顯我們的公義，來吧，我們可以在錫安報告耶和華我們神的作為。11 你們要磨尖了箭頭，抓住盾牌。耶和華定意攻擊巴比倫，將她毀滅，所以激動了瑪代君王的心，因這是耶和華報仇，就是為自己的殿報仇。 12 你們要豎立大旗攻擊巴比倫的城牆，要堅固瞭望臺，派定守望的設下埋伏。因為耶和華指著巴比倫居民所說的話、所定的意，他已經做成。 13 住在眾水之上多有財寶的啊，你的結局到了，你貪婪之量滿了！ 14 萬軍之耶和華指著自己起誓說：「我必使敵人充滿你，像螞蚱一樣，他們必呐喊攻擊你。」15 耶和華用能力創造大地，用智慧建立世界，用聰明鋪張穹蒼。 16 他一發聲，空中便有多水激動，他使雲霧從地極上騰。他造電隨雨而閃，從他府庫中帶出風來。 17 各人都成了畜類，毫無知識，各銀匠都因他的偶像羞愧。他所鑄的偶像本是虛假的，其中並無氣息， 18 都是虛無的，是迷惑人的工作，到追討的時候，必被除滅。 19 雅各的份不像這些，因他是造做萬有的主，以色列也是他產業的支派，萬軍之耶和華是他的名。20 「你是我爭戰的斧子和打仗的兵器。我要用你打碎列國，用你毀滅列邦， 21 用你打碎馬和騎馬的，用你打碎戰車和坐在其上的， 22 用你打碎男人和女人，用你打碎老年人和少年人，用你打碎壯丁和處女， 23 用你打碎牧人和他的群畜，用你打碎農夫和他一對牛，用你打碎省長和副省長。」 24 耶和華說：「我必在你們眼前，報復巴比倫人和迦勒底居民在錫安所行的諸惡。」25 耶和華說：「你這行毀滅的山哪，就是毀滅天下的山，我與你反對。我必向你伸手，將你從山巖滾下去，使你成為燒毀的山。 26 人必不從你那裡取石頭為房角石，也不取石頭為根基石，你必永遠荒涼。」這是耶和華說的。27 要在境內豎立大旗，在各國中吹角，使列國預備攻擊巴比倫，將亞拉臘、米尼、亞實基拿各國招來攻擊她。又派軍長來攻擊她，使馬匹上來如螞蚱。 28 使列國和瑪代君王，與省長和副省長，並他們所管全地之人都預備攻擊她。 29 地必震動而瘠苦，因耶和華向巴比倫所定的旨意成立了，使巴比倫之地荒涼，無人居住。 30 巴比倫的勇士止息爭戰，藏在堅壘之中，他們的勇力衰盡，好像婦女一樣。巴比倫的住處有火著起，門閂都折斷了。 31 跑報的要彼此相遇，送信的要互相迎接，報告巴比倫王說：「城的四方被攻取了， 32 渡口被占據了，葦塘被火燒了，兵丁也驚慌了。」33 萬軍之耶和華以色列的神如此說：「巴比倫城 a好像踹穀的禾場，再過片時，收割她的時候就到了。」 34 以色列人說：「巴比倫王尼布甲尼撒吞滅我，壓碎我，使我成為空虛的器皿。他像大魚將我吞下，用我的美物充滿他的肚腹，又將我趕出去。」 35 錫安的居民要說：「巴比倫以強暴待我，損害我的身體，願這罪歸給她！」耶路撒冷人要說：「願流我們血的罪歸到迦勒底的居民！」 36 所以，耶和華如此說：「我必為你申冤，為你報仇。我必使巴比倫的海枯竭，使她的泉源乾涸。 37 巴比倫必成為亂堆，為野狗的住處，令人驚駭、嗤笑，並且無人居住。 38 他們要像少壯獅子咆哮，像小獅子吼叫。 39 他們火熱的時候，我必為他們設擺酒席，使他們沉醉，好叫他們快樂，睡了長覺永不醒起。」這是耶和華說的。 40 「我必使他們像羊羔，像公綿羊和公山羊下到宰殺之地。41 「示沙克 b何竟被攻取！天下所稱讚的何竟被占據！巴比倫在列國中何竟變為荒場！ 42 海水漲起漫過巴比倫，她被許多海浪遮蓋。 43 她的城邑變為荒場、旱地、沙漠，無人居住、無人經過之地。 44 我必刑罰巴比倫的彼勒，使他吐出所吞的。萬民必不再流歸他那裡，巴比倫的城牆也必坍塌了。45 「我的民哪，你們要從其中出去，各人拯救自己，躲避耶和華的烈怒。 46 你們不要心驚膽怯，也不要因境內所聽見的風聲懼怕。因為這年有風聲傳來，那年也有風聲傳來，境內有強暴的事，官長攻擊官長。 47 \underline{日子}將到，我必刑罰巴比倫雕刻的偶像，她全地必然抱愧，她被殺的人必在其中仆倒。 48 \underline{那時}天地和其中所有的必因巴比倫歡呼，因為行毀滅的要從北方來到她那裡。」這是耶和華說的。 49 巴比倫怎樣使以色列被殺的人仆倒，照樣，她全地被殺的人也必在巴比倫仆倒。50 你們躲避刀劍的要快走，不要站住。要在遠方記念耶和華，心中追想耶路撒冷。 51 我們聽見辱罵就蒙羞，滿面慚愧，因為外邦人進入耶和華殿的聖所。 52 耶和華說：「\underline{日子}將到，我必刑罰巴比倫雕刻的偶像，通國受傷的人必唉哼。 53 巴比倫雖升到天上，雖使她堅固的高處更堅固，還有行毀滅的從我這裡到她那裡。」這是耶和華說的。54 有哀號的聲音從巴比倫出來，有大毀滅的響聲從迦勒底人之地發出。 55 因耶和華使巴比倫變為荒場，使其中的大聲滅絕。仇敵彷彿眾水，波浪砰訇，響聲已經發出。 56 這是行毀滅的臨到巴比倫，巴比倫的勇士被捉住，他們的弓折斷了，因為耶和華是施行報應的神，必定施行報應。 57 君王，名為萬軍之耶和華的說：「我必使巴比倫的首領、智慧人、省長、副省長和勇士都沉醉，使他們睡了長覺永不醒起。」 58 萬軍之耶和華如此說：「巴比倫寬闊的城牆必全然傾倒，她高大的城門必被火焚燒。眾民所勞碌的必致虛空，列國所勞碌的被火焚燒，他們都必困乏。」59 猶大王西底家在位第四年上巴比倫去的時候，瑪西雅的孫子、尼利亞的兒子西萊雅與王同去，西萊雅是王宮的大臣，先知耶利米有話吩咐他。 60 耶利米將一切要臨到巴比倫的災禍，就是論到巴比倫的一切話，寫在書上。 61 耶利米對西萊雅說：「你到了巴比倫，務要念這書上的話。 62 又說：『耶和華啊，你曾論到這地方說要剪除，甚至連人帶牲畜沒有在這裡居住的，必永遠荒涼。』




\subsubsection{以賽亞書 66:15-24}
「看哪，耶和華必在火中降臨，他的車輦像旋風，以烈怒施行報應，以火焰施行責罰。 16 因為耶和華在一切有血氣的人身上，必以火與刀施行審判，被耶和華所殺的必多。 17 那些分別為聖潔淨自己的，進入園內跟在其中一個人的後頭，吃豬肉和倉鼠並可憎之物，他們必一同滅絕。」這是耶和華說的。18 「我知道他們的行為和他們的意念。\underline{時候將到}，我必將萬民萬族聚來，看見我的榮耀。 19 我要顯神蹟在他們中間，逃脫的，我要差到列國去，就是到他施、普勒、拉弓的路德和土巴、雅完，並素來沒有聽見我名聲，沒有看見我榮耀遼遠的海島。他們必將我的榮耀傳揚在列國中。 20 他們必將你們的弟兄從列國中送回，使他們或騎馬，或坐車，坐轎，騎騾子，騎獨峰駝到我的聖山耶路撒冷，作為供物獻給耶和華，好像以色列人用潔淨的器皿盛供物奉到耶和華的殿中。」這是耶和華說的。 21 耶和華說：「我也必從他們中間取人為祭司，為利未人。」22 耶和華說：「我所要造的新天新地怎樣在我面前長存，你們的後裔和你們的名字也必照樣長存。 23 每逢月朔、安息日，凡有血氣的必來在我面前下拜。」這是耶和華說的。 24 「他們必出去觀看那些違背我人的屍首。因為他們的蟲是不死的，他們的火是不滅的。凡有血氣的，都必憎惡他們。」


\subsubsection{約珥書 2:1-21}
1 「你們要在錫安吹角，在我聖山吹出大聲！國中的居民都要發顫，因為\underline{耶和華的日子}將到，已經臨近。 2 \underline{那日}是黑暗、幽冥、密雲、烏黑的日子，好像晨光鋪滿山嶺。有一隊蝗蟲又大又強，從來沒有這樣的，以後直到萬代也必沒有。 3 牠們前面如火燒滅，後面如火焰燒盡。未到以前，地如伊甸園，過去以後，成了荒涼的曠野，沒有一樣能躲避牠們的。4 「牠們的形狀如馬，奔跑如馬兵。 5 在山頂蹦跳的響聲如車輛的響聲，又如火焰燒碎秸的響聲，好像強盛的民擺陣預備打仗。 6 牠們一來，眾民傷慟，臉都變色。 7 牠們如勇士奔跑，像戰士爬城，各都步行，不亂隊伍。 8 彼此並不擁擠，向前各行其路，直闖兵器，不偏左右。 9 牠們蹦上城，躥上牆，爬上房屋，進入窗戶如同盜賊。 10 牠們一來，地震天動，日月昏暗，星宿無光。 11 耶和華在他軍旅前發聲，他的隊伍甚大，成就他命的是強盛者。因為\underline{耶和華的日子}大而可畏，誰能當得起呢？2 「耶和華說：雖然如此，你們應當禁食、哭泣、悲哀，一心歸向我。 13 你們要撕裂心腸，不撕裂衣服，歸向耶和華你們的神，因為他有恩典，有憐憫，不輕易發怒，有豐盛的慈愛，並且後悔不降所說的災。 14 或者他轉意後悔，留下餘福，就是留下獻給耶和華你們神的素祭和奠祭，也未可知。15 「你們要在錫安吹角，分定禁食的日子，宣告嚴肅會。 16 聚集眾民，使會眾自潔，招聚老者，聚集孩童和吃奶的，使新郎出離洞房，新婦出離內室。 17 侍奉耶和華的祭司要在廊子和祭壇中間哭泣，說：『耶和華啊，求你顧惜你的百姓，不要使你的產業受羞辱，列邦管轄他們。為何容列國的人說：「他們的神在哪裡呢？」』「耶和華就為自己的地發熱心，憐恤他的百姓。 19 耶和華應允他的百姓說：我必賜給你們五穀、新酒和油，使你們飽足。我也不再使你們受列國的羞辱， 20 卻要使北方來的軍隊遠離你們，將他們趕到乾旱荒廢之地，前隊趕入東海，後隊趕入西海。因為他們所行的大惡 b，臭氣上升，腥味騰空。」21 地土啊，不要懼怕，要歡喜快樂！因為耶和華行了大事。 22 田野的走獸啊，不要懼怕！因為曠野的草發生，樹木結果，無花果樹、葡萄樹也都效力。 23 錫安的民哪，你們要快樂，為耶和華你們的神歡喜！因他賜給你們合宜的秋雨，為你們降下甘霖，就是秋雨、春雨，和先前一樣。 24 禾場必滿了麥子，酒榨與油榨必有新酒和油盈溢。 25 「我打發到你們中間的大軍隊，就是蝗蟲、蝻子、螞蚱、剪蟲，那些年所吃的，我要補還你們。 26 你們必多吃而得飽足，就讚美為你們行奇妙事之耶和華你們神的名。我的百姓必永遠不致羞愧。 27 你們必知道我是在以色列中間，又知道我是耶和華你們的神，在我以外並無別神。我的百姓必永遠不致羞愧。28 「以後，我要將我的靈澆灌凡有血氣的，你們的兒女要說預言，你們的老年人要做異夢，少年人要見異象。 29 在\underline{那些日子}，我要將我的靈澆灌我的僕人和使女。 30 在天上地下，我要顯出奇事，有血，有火，有煙柱。 31 日頭要變為黑暗，月亮要變為血，這都在\underline{耶和華大而可畏的日子}未到以前。32「\underline{到那時候}，凡求告耶和華名的，就必得救。因為照耶和華所說的，在錫安山耶路撒冷必有逃脫的人，在剩下的人中必有耶和華所召的。


\subsubsection{約珥書 3:1-21}
1 「到\underline{那日}，我使猶大和耶路撒冷被擄之人歸回的時候， 2 我要聚集萬民，帶他們下到約沙法谷，在那裡施行審判。因為他們將我的百姓，就是我的產業以色列，分散在列國中，又分取我的地土， 3 且為我的百姓拈鬮，將童子換妓女，賣童女買酒喝。 4 推羅、西頓和非利士四境的人哪，你們與我何干？你們要報復我嗎？若報復我，我必使報應速速歸到你們的頭上。 5 你們既然奪取我的金銀，又將我可愛的寶物帶入你們宮殿 a， 6 並將猶大人和耶路撒冷人賣給希臘人，使他們遠離自己的境界， 7 我必激動他們離開你們所賣到之地，又必使報應歸到你們的頭上。 8 我必將你們的兒女賣在猶大人的手中，他們必賣給遠方示巴國的人。這是耶和華說的。9 「當在萬民中宣告說：要預備打仗，激動勇士，使一切戰士上前來！ 10 要將犁頭打成刀劍，將鐮刀打成戈矛。軟弱的要說：『我有勇力！』」 11 四圍的列國啊，你們要速速地來，一同聚集！耶和華啊，求你使你的大能者降臨！ 12 「萬民都當興起，上到約沙法谷，因為我必坐在那裡，審判四圍的列國。 13 開鐮吧！因為莊稼熟了。踐踏吧！因為酒榨滿了，酒池盈溢，他們的罪惡甚大。14 「許多許多的人在斷定谷，因為耶和華的日子臨近斷定谷。 15 日月昏暗，星宿無光。 16 耶和華必從錫安吼叫，從耶路撒冷發聲，天地就震動。耶和華卻要做他百姓的避難所，做以色列人的保障。 17 你們就知道我是耶和華你們的神，且又住在錫安我的聖山。那時，耶路撒冷必成為聖，外邦人不再從其中經過。18 「到\underline{那日}，大山要滴甜酒，小山要流奶子，猶大溪河都有水流，必有泉源從耶和華的殿中流出來，滋潤什亭谷。 19 埃及必然荒涼，以東變為淒涼的曠野，都因向猶大人所行的強暴，又因在本地流無辜人的血。「但猶大必存到永遠，耶路撒冷必存到萬代。 21 我未曾報復流血的罪，現在我要報復，因為耶和華住在錫安。」

\subsubsection{撒迦利亞書 14:1-21}

1 \underline{耶和華的日子}臨近，你的財物必被搶掠，在你中間分散。 2 「因為我必聚集萬國與耶路撒冷爭戰，城必被攻取，房屋被搶奪，婦女被玷汙，城中的民一半被擄去，剩下的民仍在城中，不致剪除。」 3 那時耶和華必出去與那些國爭戰，好像從前爭戰一樣。4 \underline{那日}，他的腳必站在耶路撒冷前面朝東的橄欖山上。這山必從中間分裂，自東至西成為極大的谷，山的一半向北挪移，一半向南挪移。 5 「你們要從我山的谷中逃跑，因為山谷必延到亞薩。你們逃跑，必如猶大王烏西雅年間的人逃避大地震一樣。」耶和華我的神必降臨，有一切聖者同來。 6 \underline{那日}必沒有光，三光必退縮。 7 \underline{那日}必是耶和華所知道的，不是白晝，也不是黑夜，到了晚上才有光明。 8 \underline{那日}必有活水從耶路撒冷出來，一半往東海流，一半往西海流，冬夏都是如此。
9 耶和華必做全地的王，那日耶和華必為獨一無二的，他的名也是獨一無二的。 10 全地，從迦巴直到耶路撒冷南方的臨門，要變為亞拉巴。耶路撒冷必仍居高位，就是從便雅憫門到第一門之處，又到角門，並從哈楠業樓直到王的酒榨。 11 人必住在其中，不再有咒詛，耶路撒冷人必安然居住。12 耶和華用災殃攻擊那與耶路撒冷爭戰的列國人，必是這樣：他們兩腳站立的時候，肉必消沒，眼在眶中乾癟，舌在口中潰爛。 13 \underline{那日}，耶和華必使他們大大擾亂，他們各人彼此揪住，舉手攻擊。 14 猶大也必在耶路撒冷爭戰。那時四圍各國的財物，就是許多金銀、衣服，必被收聚。 15 那臨到馬匹、騾子、駱駝、驢和營中一切牲畜的災殃是與那災殃一般。16 所有來攻擊耶路撒冷列國中剩下的人，必年年上來敬拜大君王萬軍之耶和華，並守住棚節。 17 地上萬族中，凡不上耶路撒冷敬拜大君王萬軍之耶和華的，必無雨降在他們的地上。 18 埃及族若不上來，雨也不降在他們的地上。凡不上來守住棚節的列國人，耶和華也必用這災攻擊他們。 19 這就是埃及的刑罰和那不上來守住棚節之列國的刑罰。 20 當\underline{那日}，馬的鈴鐺上必有「歸耶和華為聖的」這句話。耶和華殿內的鍋必如祭壇前的碗一樣。 21 凡耶路撒冷和猶大的鍋都必歸萬軍之耶和華為聖，凡獻祭的都必來取這鍋，煮肉在其中。當\underline{那日}，在萬軍之耶和華的殿中必不再有迦南人。


\subsubsection{瑪拉基書 3:1-4，16-18}
萬軍之耶和華說：「我要差遣我的使者在我前面預備道路。你們所尋求的主必忽然進入他的殿，立約的使者，就是你們所仰慕的，快要來到。「他來的日子，誰能當得起呢？他顯現的時候，誰能立得住呢？因為他如煉金之人的火，如漂布之人的鹼。 3 他必坐下如煉淨銀子的，必潔淨利未人，熬煉他們像金銀一樣，他們就憑公義獻供物給耶和華。 4 \underline{那時}猶大和耶路撒冷所獻的供物必蒙耶和華悅納，彷彿古時之日，上古之年。」16  那時，敬畏耶和華的彼此談論，耶和華側耳而聽，且有紀念冊在他面前，記錄那敬畏耶和華思念他名的人。 17 萬軍之耶和華說：「在我\underline{所定的日子}他們必屬我，特特歸我，我必憐恤他們，如同人憐恤服侍自己的兒子。 18 那時你們必歸回，將善人和惡人，侍奉神的和不侍奉神的，分別出來。

\subsubsection{瑪拉基書 4:1-5}
1 萬軍之耶和華說：「\underline{那日}臨近，勢如燒著的火爐，凡狂傲的和行惡的必如碎秸，在\underline{那日}必被燒盡，根本、枝條一無存留。 2 但向你們敬畏我名的人，必有公義的日頭出現，其光線有醫治之能。你們必出來，跳躍如圈裡的肥犢。 3 你們必踐踏惡人，在我\underline{所定的日子}，他們必如灰塵在你們腳掌之下。這是萬軍之耶和華說的。4 「你們當記念我僕人摩西的律法，就是我在何烈山為以色列眾人所吩咐他的律例、典章。「看哪，\underline{耶和華大而可畏之日}未到以前，我必差遣先知以利亞到你們那裡去。 6 他必使父親的心轉向兒女，兒女的心轉向父親，免得我來咒詛遍地。」


\subsubsection{使徒行傳 2:14-22}
14 彼得和十一個使徒站起，高聲說：「猶太人和一切住在耶路撒冷的人哪，這件事你們當知道，也當側耳聽我的話。 15 你們想這些人是醉了，其實不是醉了，因為時候剛到巳初。 16 這正是先知約珥所說的： 17 『神說：在\underline{末後的日子}，我要將我的靈澆灌凡有血氣的，你們的兒女要說預言，你們的少年人要見異象，老年人要做異夢。 18 在\underline{那些日子}，我要將我的靈澆灌我的僕人和使女，他們就要說預言。 19 在天上我要顯出奇事，在地下我要顯出神蹟，有血，有火，有煙霧。 20 日頭要變為黑暗，月亮要變為血，這都在\underline{主大而明顯的日子}未到以前。 21 到那時候，凡求告主名的，就必得救。』 22 以色列人哪，請聽我的話：神藉著拿撒勒人耶穌在你們中間施行異能、奇事、神蹟，將他證明出來，


\subsubsection{帖撒羅尼迦後書 1:3-10}
3 弟兄們，我們該為你們常常感謝神，這本是合宜的，因你們的信心格外增長，並且你們眾人彼此相愛的心也都充足。 4 甚至我們在神的各教會裡為你們誇口，都因你們在所受的一切逼迫患難中，仍舊存忍耐和信心。 5 這正是神公義判斷的明證，叫你們可算配得神的國，你們就是為這國受苦。 6 神既是公義的，就必將患難報應那加患難給你們的人， 7 也必使你們這受患難的人與我們同得平安。那時，主耶穌同他有能力的天使從天上在火焰中顯現， 8 要報應那不認識神和那不聽從我主耶穌福音的人。 9 他們要受刑罰，就是永遠沉淪，離開主的面和他權能的榮光。 10 這正是主降臨，要在他聖徒的身上得榮耀，又在一切信的人身上顯為稀奇的\underline{那日子}。我們對你們作的見證，你們也信了。


\subsubsection{彼得後書 3:3-13}
 3 第一要緊的，該知道在末世必有好譏誚的人隨從自己的私慾出來譏誚說： 4 「主要降臨的應許在哪裡呢？因為從列祖睡了以來，萬物與起初創造的時候仍是一樣。」 5 他們故意忘記：從太古憑神的命有了天，並從水而出、藉水而成的地； 6 故此，當時的世界被水淹沒就消滅了。 7 但現在的天地還是憑著那命存留，直留到不敬虔之人受審判遭沉淪的日子，用火焚燒。8 親愛的弟兄啊，有一件事你們不可忘記，就是主看一日如千年，千年如一日。 9 主所應許的尚未成就，有人以為他是耽延，其實不是耽延，乃是寬容你們，不願有一人沉淪，乃願人人都悔改。 10 但主的日子要像賊來到一樣。那日，天必大有響聲廢去，有形質的都要被烈火銷化，地和其上的物都要燒盡了。 11 這一切既然都要如此銷化，你們為人該當怎樣聖潔，怎樣敬虔， 12 切切仰望神的日子來到！在\underline{那日}，天被火燒就銷化了，有形質的都要被烈火熔化。 13 但我們照他的應許，盼望新天新地，有義居在其中。






\subsection{七靈,七個金燈臺,七星,七眼}
\subsubsection{出埃及記 25:37，37:23}
\begin{itemize} 
\itemsep-.25em
\item 要做燈臺的七個燈盞，祭司要點這燈，使燈光對照(出 25:37)。
\item 用精金做燈臺的七個燈盞，並燈臺的蠟剪和蠟花盤(出 37:23)。
%\item 我說：「我看見了一個純金的燈臺，頂上有燈盞，燈臺上有七盞燈，每盞有七個管子(撒迦利亞書 4:2)。
\item 我又看見寶座與四活物並長老之中有羔羊站立，像是被殺過的，有七角七眼，就是神的七靈，奉差遣往普天下去的(啟示錄 5:6)。
\end{itemize} 

\subsubsection{撒迦利亞書 3:8-10，4:6-10}
\begin{itemize} 
\itemsep-.25em
\item 「大祭司約書亞啊，你和坐在你面前的同伴都當聽（他們是做預兆的）：我必使我僕人大衛的苗裔發出。 9 看哪，我在約書亞面前所立的石頭，在一塊石頭上有七眼。萬軍之耶和華說：我要親自雕刻這石頭，並要在一日之間除掉這地的罪孽。 10 當那日，你們各人要請鄰舍坐在葡萄樹和無花果樹下。這是萬軍之耶和華說的。」
\item 他對我說：「這是耶和華指示所羅巴伯的。萬軍之耶和華說：不是倚靠勢力，不是倚靠才能，乃是倚靠我的靈方能成事。 7 大山哪，你算什麼呢？在所羅巴伯面前，你必成為平地。他必搬出一塊石頭，安在殿頂上，人且大聲歡呼說：『願恩惠，恩惠歸於這殿！』」 8 耶和華的話又臨到我說： 9 「所羅巴伯的手立了這殿的根基，他的手也必完成這工，你就知道萬軍之耶和華差遣我到你們這裡來了。 10 誰藐視這日的事為小呢？這七眼乃是耶和華的眼睛，遍察全地，見所羅巴伯手拿線砣就歡喜。」
\end{itemize} 

\subsubsection{啟示錄 1:4-6，12-20，2:1-3，5，3:1-2， 5:6}
\begin{itemize} 
\itemsep-.25em
\item 約翰寫信給亞細亞的七個教會。但願從那昔在、今在、以後永在的神和他寶座前的七靈， 5 並那誠實作見證的、從死裡首先復活、為世上君王元首的耶穌基督，有恩惠、平安歸於你們！他愛我們，用自己的血使我們脫離罪惡， 6 又使我們成為國民，做他父神的祭司。但願榮耀、權能歸給他，直到永永遠遠！阿們。
\item 我轉過身來，要看是誰發聲與我說話；既轉過來，就看見七個金燈臺， 13 燈臺中間有一位好像人子，身穿長衣，直垂到腳，胸間束著金帶。 14 他的頭與髮皆白，如白羊毛，如雪，眼目如同火焰， 15 腳好像在爐中鍛煉光明的銅，聲音如同眾水的聲音。 16 他右手拿著七星，從他口中出來一把兩刃的利劍，面貌如同烈日放光。17 我一看見，就仆倒在他腳前，像死了一樣。他用右手按著我說：「不要懼怕。我是首先的，我是末後的， 18 又是那存活的；我曾死過，現在又活了，直活到永永遠遠，並且拿著死亡和陰間的鑰匙。 19 所以你要把所看見的和現在的事，並將來必成的事，都寫出來。 20 論到你所看見在我右手中的七星和七個金燈臺的奧祕，那七星就是七個教會的使者，七燈臺就是七個教會。
\item 1 「你要寫信給以弗所教會的使者說：『那右手拿著七星、在七個金燈臺中間行走的說： 2 我知道你的行為、勞碌、忍耐，也知道你不能容忍惡人。你也曾試驗那自稱為使徒卻不是使徒的，看出他們是假的來。 3 你也能忍耐，曾為我的名勞苦，並不乏倦。5 所以，應當回想你是從哪裡墜落的，並要悔改，行起初所行的事。你若不悔改，我就臨到你那裡，把你的燈臺從原處挪去。
\item 「你要寫信給撒狄教會的使者說：『那有神的七靈和七星的說：我知道你的行為，按名你是活的，其實是死的。 2 你要警醒，堅固那剩下將要衰微的，因我見你的行為在我神面前，沒有一樣是完全的。
\item 我又看見寶座與四活物並長老之中有羔羊站立，像是被殺過的，有七角七眼，就是神的七靈，奉差遣往普天下去的。
\end{itemize} 

















